\chapter{Marked duality and coefficient polarity}
\label{v4-arith:md:chapter}

Arithmetic duality must transport an operation before it can explain a
functional equation. Its first datum is a pairing between specified
objects. The pairing determines the direction of a transpose and the
tensor product in which evaluation is defined. A topology then decides
which transposes remain continuous. These choices already distinguish
several dualities on the polynomial oscillator representation.

We begin with finite coefficients, where evaluation is explicit. We
then pass to compact modules, marked derived modules and coalgebras.
The completed oscillator shows why a continuous coaction can exist
without a continuous evaluation on the proposed tensor. Full collision
coefficients admit a different duality, whose field realization can be
constructed on a marked charged-current domain. The arithmetic question
is to compare such operations with theta and Mellin reflection while
retaining the vacuum and both polar terms.

\section{The pairing fixes the dual}
\label{v4-arith:md:sec:finite}

Let $R$ be a commutative ring. For a finite projective $R$-module $P$,
put $P^\vee=\operatorname{Hom}_R(P,R)$. Evaluation is
$P^\vee\otimes_RP\to R$, $f\otimes p\mapsto f(p)$.

A map $a:P\to Q$ has transpose
$a^\vee:Q^\vee\to P^\vee$, $f\mapsto f\circ a$. Thus
\begin{equation}\label{v4-arith:md:eq:transpose}
 \langle a(p),f\rangle=\langle p,a^\vee(f)\rangle,
 \qquad (ba)^\vee=a^\vee b^\vee.
\end{equation}
These formulas determine variance. They do not assert that a given
bilinear form identifies $P$ with its dual.

\begin{proposition}[Finite projective polarity]
\label{v4-arith:md:prop:finite}
Evaluation gives natural isomorphisms
\[
 P\xrightarrow{\sim}P^{\vee\vee},\qquad
 Q^\vee\otimes_RP^\vee\xrightarrow{\sim}(P\otimes_RQ)^\vee.
\]
The second isomorphism uses evaluation in the order of the displayed
factors. For a ring map $R\to R'$, the natural map
\[
 P^\vee\otimes_RR'\longrightarrow
 \operatorname{Hom}_{R'}(P\otimes_RR',R')
\]
is an isomorphism. These statements extend to perfect complexes, with
derived tensors and the cohomological Koszul signs.
\end{proposition}

\begin{proof}
For $P=R^n$ the evaluation matrix in a basis and its dual is the
identity. The tensor assertion and coefficient comparison are the
corresponding identities for finite matrices. A finite projective
module is a direct summand of a finite free module, so the identities
restrict to it.

Apply the same argument to bounded complexes of finite
projectives. For a degree-$j$ functional the dual differential is
$d(f)=(-1)^{j+1}f\,d_P$. This makes evaluation a chain map.
Finite cones and direct summands preserve all three comparisons.
\end{proof}

If $C$ is a finite projective coalgebra, its coproduct transposes to
an associative multiplication on $C^\vee$. The coalgebra counit
transposes to the algebra unit. Conversely, an algebra multiplication
transposes to a coproduct on its finite projective dual. Each identity
is obtained by applying the same duality to its defining diagram.
This is coefficient transposition. The cobar algebra of $C$ instead
uses tensor words in its shifted coaugmentation quotient and a
differential determined by the reduced coproduct.

\subsection{An integral divided-power example}

Consider $R[x]$ with $x$ primitive:
$\Delta x=x\otimes1+1\otimes x$. Give $x$ weight one, and
write $f_n$ for the coefficient functional on $x^n$.
The \emph{restricted graded dual} is
$\bigoplus_{n\geq0}Rf_n$, the direct sum of the dual weight spaces.

Transposing the coproduct gives
\begin{equation}\label{v4-arith:md:eq:divided}
 f_af_b=\binom{a+b}{a}f_{a+b},\qquad
 f_1^n=n!f_n,\qquad f_0=1.
\end{equation}
Indeed the coefficient of $x^a\otimes x^b$ in
$\Delta(x^{a+b})$ is the binomial coefficient. The identity
\[
 \binom{a+b}{a}\binom{a+b+c}{a+b}
 =\frac{(a+b+c)!}{a!b!c!}
 =\binom{b+c}{b}\binom{a+b+c}{a}
\]
proves associativity over the integers, and hence over every $R$.
It involves no division in $R$.

Over a field of characteristic zero the substitution
$f_n\mapsto y^n/n!$ identifies this algebra with a polynomial
algebra. Over $\mathbb F_p$, however, $f_1^p=0$ while $f_p\ne0$.
The integral dual therefore retains divided powers. For several
primitive variables the formula is the product of the corresponding
binomial coefficients. The energy grading of the oscillator makes
each weight finite, so the same calculation applies weight by weight.

The full linear dual of $\bigoplus_nRx^n$ is instead
$\prod_nRf_n$. Its multiplication still has finite sums in each
output weight. Whether its evaluation extends to a tensor completion
is a separate question, considered in
Section~\ref{v4-arith:md:sec:completed-modules}.

\section{Compact coefficients and their discrete duals}
\label{v4-arith:md:sec:compact}

For a locally compact Hausdorff abelian group $G$, its Pontryagin dual
is $G^*=\operatorname{Hom}_{\mathrm{cts}}(G,\mathbb R/\mathbb Z)$,
with the compact-open topology. The classical duality theorem says
that evaluation $G\to G^{**}$ is a topological isomorphism. It
exchanges compact and discrete groups. See Roeder,
\emph{Category theory applied to Pontryagin duality}, Theorems 2, 4
and 5. The specialization needed for Iwasawa coefficients has an
elementary proof through finite groups.

Fix a prime $p$ and put $D_p=\mathbb Q_p/\mathbb Z_p$, with the
discrete topology. A \emph{profinite $\mathbb Z_p$-module} is an
inverse limit of finite $p$-primary modules with its inverse-limit
topology. For such an $M$, and for a discrete $p$-primary torsion
module $N$, define
\[
 M^\vee=\operatorname{Hom}_{\mathrm{cts}}(M,D_p),\qquad
 N^\vee=\operatorname{Hom}_{\mathbb Z_p}(N,D_p).
\]
The first dual is discrete. Give the second the inverse-limit
topology obtained by restricting characters to finite submodules
of $N$. Every finitely generated submodule of $N$ is finite.

\begin{theorem}[Compact and discrete coefficient duality]
\label{v4-arith:md:thm:pontryagin}
These duals are exact contravariant inverse equivalences between
profinite $\mathbb Z_p$-modules and discrete $p$-primary torsion
modules. Their double-dual comparisons are evaluation maps.
\end{theorem}

\begin{proof}
For $\mathbb Z/p^n$, a character is determined by the image of $1$,
an element of $D_p[p^n]$. The pairing of these two cyclic groups is
perfect. Finite direct sums prove the assertion for every finite
$p$-group.

A continuous character on $M$ has an open kernel, so it factors
through a finite quotient. Consequently
\[
 M^\vee=\varinjlim_U(M/U)^\vee,
 \qquad N^\vee=\varprojlim_{N_0\subset N}N_0^\vee,
\]
where $U$ runs through open submodules and $N_0$ through finite
submodules. Taking a second dual and using the finite evaluation
identities recovers $M$ with its topology and $N$ with its discrete
topology. The comparisons are natural.

To prove exactness, let $L\subset M$ be a closed submodule and let
$\chi:L\to D_p$ be continuous. Choose an open submodule $U$ of
$M$ such that $L\cap U\subset\ker\chi$. The character of
$L/(L\cap U)\subset M/U$ extends to $M/U$, since $D_p$ is
divisible and hence injective as a $\mathbb Z$-module. Pulling back
gives a continuous extension to $M$. The kernel of restriction is
exactly the dual of $M/L$. Double duality gives exactness on the
discrete side as well.
\end{proof}

The same $D_p$ describes circle characters of a profinite $p$-group.
Choose an arc around zero in $\mathbb R/\mathbb Z$ containing no
nonzero subgroup. The inverse image of this arc contains an open
subgroup $U$.

Its image must be zero. Every character therefore
factors through a finite $p$-group and has $p$-primary torsion image.
Here $D_p$ is identified with the $p$-power torsion subgroup of the
circle. It is given its discrete topology in the preceding duality.

\subsection{The dual tensor and the Iwasawa algebra}

For profinite modules define
\[
 M\widehat\otimes_{\mathbb Z_p}N
 =\varprojlim_{U,V}(M/U)\otimes_{\mathbb Z_p}(N/V).
\]
For discrete $p$-primary modules use the transported product
\begin{equation}\label{v4-arith:md:eq:dual-tensor}
 A\odot B=(A^\vee\widehat\otimes_{\mathbb Z_p}B^\vee)^\vee.
\end{equation}
Its unit is $D_p$, since $D_p^\vee=\mathbb Z_p$.
Associativity and symmetry are transported through the natural
double-dual maps. Thus a compact unital algebra $\Lambda$ has a
dual coalgebra $C=\Lambda^\vee$ for $\odot$, with
$\Delta:C\to C\odot C$ and counit $C\to D_p$.

Dualizing an action $M\widehat\otimes\Lambda\to M$ gives a
coaction $M^\vee\to M^\vee\odot C$. Its associativity diagram
becomes the coassociativity diagram. Dualizing again recovers the
action. Hence continuous compact right $\Lambda$-modules and
discrete right $C$-comodules for $\odot$ are exactly
anti-equivalent. Exactness extends this comparison to bounded
derived categories, with the dual differential used above.

The transported tensor is essential. Ordinary tensor gives
$D_p\otimes_{\mathbb Z_p}D_p=0$: if $p^na=0$, write $b=p^nc$
and compute $a\otimes b=p^na\otimes c=0$.
Moreover $\operatorname{Hom}_{\mathbb Z_p}(D_p,\mathbb Z_p)=0$.
Thus $D_p$ cannot be a nonzero ordinary counital
$\mathbb Z_p$-coalgebra. Its counit for $\odot$ takes values in
the correct unit $D_p$.

Let $G$ be a profinite group and put
\[
 \Lambda(G)=\mathbb Z_p[[G]]
 =\varprojlim_{U,n}(\mathbb Z/p^n)[G/U],
\]
where $U$ is open and normal. The finite coefficient pairing
identifies the dual of $(\mathbb Z/p^n)[G/U]$ with the functions
$G/U\to D_p[p^n]$. Therefore
\begin{equation}\label{v4-arith:md:eq:iwasawa-dual}
 \Lambda(G)^\vee=C_{\mathrm{lc}}(G,D_p).
\end{equation}
Every locally constant function on the compact group has finite
image and factors through some finite quotient. Its values are
killed by a common power of $p$. This proves surjectivity of the
displayed identification as well as injectivity.

For left translation on $\Lambda(G)$, choose the contragredient
left action on its dual. It is
$(g\cdot f)(h)=f(g^{-1}h)$. This follows by evaluating a character
on the group-ring basis element $[h]$. The full local abelian Galois
group can be used for $G$, but it is not the chosen procyclic group
whose Iwasawa algebra is $\mathbb Z_p[[T]]$.

There is also a natural identity
\begin{equation}\label{v4-arith:md:eq:reduction-torsion}
 (M/pM)^\vee=M^\vee[p].
\end{equation}
Indeed a character vanishes on $pM$ exactly when multiplication by
$p$ kills it. This is an identity of ordinary modules. Derived
reduction retains additional Tor groups.

\subsection{Matlis duality and the coefficient object}

For a complete Noetherian local ring $(R,\mathfrak m,k)$, let $E_R$
be an injective hull of $k$. Thus $E_R$ is injective and contains
$k$ as an essential submodule: every nonzero submodule meets $k$.
Define the \emph{Matlis dual} by
$D_R(M)=\operatorname{Hom}_R(M,E_R)$.
Its coefficient object is $E_R$, rather than $R$.

\begin{theorem}[Matlis duality]
\label{v4-arith:md:thm:matlis}
The functor $D_R$ is an exact anti-equivalence between finite
$R$-modules and Artinian $R$-modules. Evaluation identifies either
kind of module with its double dual.
\end{theorem}

\begin{proof}
The injectivity of $E_R$ makes $D_R$ exact. The structural facts
are that $E_R$ is Artinian, that $\operatorname{End}_R(E_R)=R$,
and that every Artinian $M$ has a finite copresentation
$0\to M\to E_R^a\to E_R^b$. They are proved in the Stacks Project,
Tags 08Z6--08Z8. A finite module has a presentation
$R^b\to R^a\to M\to0$.

Dualizing twice and using
$D_R(R)=E_R$, $D_R(E_R)=R$ proves biduality by kernels and
cokernels. The same argument applied to the copresentation proves
biduality on the Artinian side. Naturality of evaluation proves
that the two functors are inverse. See also Tag 08Z9.
\end{proof}

Completeness enters the identity for $\operatorname{End}_R(E_R)$.
Without it that endomorphism ring is $\widehat R$. Indeed
$E_R=\bigcup_nE_R[\mathfrak m^n]$, and endomorphisms restrict to
the finite-length duality on each quotient. Their inverse limit is
$\varprojlim_nR/\mathfrak m^n$.
For $R=\mathbb Z_p$, the injective hull is $D_p$, so Matlis duality
agrees with the preceding pairing on finite $\mathbb Z_p$-modules.
Its stated domain is smaller: $\bigoplus_{n\geq0}\mathbb F_p$
is discrete $p$-primary torsion but is not Artinian.

Replacing $E_R$ by $R$ can lose the entire module. For example,
$\operatorname{Hom}_{\mathbb Z_p}(\mathbb F_p,\mathbb Z_p)=0$,
whereas $D_{\mathbb Z_p}(\mathbb F_p)\cong\mathbb F_p$.
A claim of biduality must therefore specify its dualizing object
and its module class.

\section{Marked modules and balanced interfaces}
\label{v4-arith:md:sec:marked}

The dual of a module can retain a chosen boundary object rather
than a scalar coefficient ring. The domain of biduality is then
determined by the objects generated from that boundary.

Let $A$ be a dg algebra and $P$ a right $A$-module. Choose a
derived model for
\[
 E_P=\operatorname{RHom}_A(P,P),
\]
with multiplication given by composition. Then $P$ is an
$(E_P,A)$-bimodule. Write $\operatorname{thick}_A(P)$ for the
smallest full stable subcategory of $D(A)$ containing $P$ and
closed under direct summands. Its objects are obtained by finitely
many shifts, cones and summands. A left $E_P$-module is
\emph{perfect} when it lies in $\operatorname{thick}_{E_P}(E_P)$.

\begin{proposition}[Duality generated by a marked module]
\label{v4-arith:md:prop:thick}
The functors
\[
 F_P=\operatorname{RHom}_A(-,P),\qquad
 G_P=\operatorname{RHom}_{E_P}(-,P)
\]
give inverse anti-equivalences
$\operatorname{thick}_A(P)^{\mathrm{op}}\simeq
\operatorname{Perf}(E_P)$. The first Hom has a left $E_P$-action
by postcomposition. The second has a right $A$-action through $P$.
\end{proposition}

\begin{proof}
Both functors are exact and preserve finite direct summands.
They send $P$ to the regular module $E_P$ and $E_P$ to $P$.
The evaluation maps from an object to its double Hom are
equivalences on these generators. The class where either
evaluation is an equivalence is closed under shifts, cones and
summands. Hence both evaluations are equivalences on the stated
categories. The evaluation identities give the inverse functors
on mapping complexes and their compositions.
\end{proof}

Now let $S$ be a Noetherian ring of finite global dimension and
let $A$ be an ordinary augmented $S$-algebra, finite over $S$.
Assume that its augmentation ideal $I$ is nilpotent. A complex
of right $A$-modules is \emph{coefficient-perfect} when its
restriction to $S$ is perfect. Put
$E=\operatorname{RHom}_A(S,S)$.

\begin{theorem}[The coefficient-perfect domain]
\label{v4-arith:md:thm:marked}
Coefficient-perfect $A$-modules form $\operatorname{thick}_A(S)$.
The functor $F_A(M)=\operatorname{RHom}_A(M,S)$ therefore gives
an anti-equivalence from this category to perfect left
$E$-modules. It sends $S$ to $E$ and $A$ to the augmentation
module $S$ of $E$.
\end{theorem}

\begin{proof}
The object $S$ is coefficient-perfect, and this property is
preserved by finite cones and summands. This proves one inclusion.
Conversely, a coefficient-perfect complex has bounded, finite
$S$-module cohomology. A finite Postnikov filtration reduces the
claim to each cohomology module $M$. Filter $M$ by the powers of
$I$.

Since $I$ is nilpotent, there are finitely many quotients,
each a finite $S$-module with $A$-action through the augmentation.
Every such quotient has a finite projective $S$-resolution, so
belongs to $\operatorname{thick}_A(S)$. Finite extensions prove
the other inclusion. Apply Proposition~\ref{v4-arith:md:prop:thick}
with $P=S$. The boundary values follow from the two displayed
derived Hom complexes.
\end{proof}

Let $\mathcal C=\operatorname{thick}_A(S)$. Its pro-category
retains an inverse system as an object, including its transition
maps. The extension of the preceding duality is
\begin{equation}\label{v4-arith:md:eq:pro-ind}
 \operatorname{Pro}(\mathcal C)^{\mathrm{op}}
 \simeq\operatorname{Ind}(\operatorname{Perf}(E))
 \simeq D(E),\qquad
 \{M_i\}\longmapsto\varinjlim_i\operatorname{RHom}_A(M_i,S).
\end{equation}
The direct limit is taken over the opposite of the cofiltered
index category.

To see the first equivalence, apply the small
anti-equivalence to the formula for maps of pro-objects,
$\varprojlim_j\varinjlim_i\operatorname{Map}(M_i,N_j)$.
The resulting formula is precisely that for the corresponding
ind-objects, with the arguments reversed. The second equivalence
uses the regular module $E$ as a compact generator of $D(E)$.
Taking the inverse limit of $\{M_i\}$ inside $D(A)$ before
dualizing is a different operation.

\subsection{Two tests for the domain}

For $A=k[\epsilon]/(\epsilon^2)$, with $k$ a field, the
augmentation has the free resolution
\[
 \cdots\xrightarrow{\epsilon}A
 \xrightarrow{\epsilon}A\xrightarrow{\epsilon}A\longrightarrow k.
\]
The kernel and image of each multiplication map are both
$k\epsilon$, which proves exactness. Applying
$\operatorname{Hom}_A(-,k)$ gives zero differentials and one
copy of $k$ in every nonnegative degree.

The periodic extension
class generates their Yoneda products: splicing its first
$n$ steps gives the nonzero class in degree $n$. Consequently
\begin{equation}\label{v4-arith:md:eq:dual-numbers}
 E\simeq k\langle t\rangle,
 \qquad |t|=1,\qquad t^2\ne0.
\end{equation}
To obtain the dg model, lift the degree-one class to the
endomorphism complex of the resolution. Its components are
alternating scalar identities, so it anticommutes with the
differential.

Its powers represent the nonzero extension
classes just computed. The resulting map from the free
associative dg algebra $k\langle t\rangle$ is a quasi-isomorphism.
This algebra is not required to be graded commutative.
The nonzero Ext groups in arbitrarily high degrees show that $k$
is coefficient-perfect and is not perfect over $A$.

The same small duality sends $\operatorname{Perf}(A)$ to
$\operatorname{thick}_E(k)$. Its ind-extension generates
$\operatorname{Loc}_E(k)$, the full stable subcategory closed
under colimits and generated by $k$.
Inverting $t$ annihilates $k$, and hence this whole localizing
subcategory, but does not annihilate $E$. Thus the
algebra-perfect and coefficient-perfect extensions differ.
In this example $k$ is perfect over $E$, as the two-cell
resolution given by multiplication by $t$ shows.

Nilpotence of the augmentation ideal also matters. Take $A=k[x]$
and $E=\Lambda_k(\eta)$, with $|\eta|=1$ and zero differential.
The augmentation $E\to k$ has a semifree resolution with
generators $e_n$, $n\geq0$, of degree zero and
\[
 de_0=0,\qquad de_n=e_{n-1}\eta\quad(n\geq1).
\]
Every element is a finite sum of generators with coefficients
in $E$.

All positive-degree cycles are boundaries, and its
degree-zero cohomology is generated by $e_0$.
Applying $\operatorname{Hom}_E(-,k)$ gives $\prod_{n\geq0}k$
in degree zero. The endomorphism $e_n\mapsto e_{n-1}$, with
$e_{-1}=0$, identifies composition with multiplication of
formal power series. Hence
\begin{equation}\label{v4-arith:md:eq:polynomial-double}
 \operatorname{RHom}_A(A,k)=k,
 \qquad\operatorname{RHom}_E(k,k)=k[[x]].
\end{equation}

Each formal series acts by a finite sum on every $e_n$.
The second dual has completed $k[x]$. It has not returned the
polynomial algebra. Retaining the tower $\{A/(x^n)\}_n$
records the missing topology.

\subsection{The identity interface}

Choose right modules $P_A,P_B,P_C$ over dg algebras $A,B,C$,
and let $E_A=\operatorname{RHom}_A(P_A,P_A)$, with analogous
notation for $B,C$. An $(A,B)$-bimodule $K$ acts on right
modules by $M\mapsto M\otimes_A^{\mathbb L}K$.
Assume
\[
 P_A\otimes_A^{\mathbb L}K\in\operatorname{thick}_B(P_B),
 \qquad P_B\otimes_B^{\mathbb L}H\in\operatorname{thick}_C(P_C)
\]
for a second bimodule ${}_BH_C$. Define the transformed interfaces
\[
 {}_{E_B}L_K{}_{E_A}
 =\operatorname{RHom}_B(P_A\otimes_A^{\mathbb L}K,P_B),
 \qquad
 {}_{E_C}L_H{}_{E_B}
 =\operatorname{RHom}_C(P_B\otimes_B^{\mathbb L}H,P_C).
\]
Their left actions are postcomposition and their right actions
are precomposition. These choices fix the order of the product.

\begin{proposition}[Composition of marked interfaces]
\label{v4-arith:md:prop:interfaces}
Composition induces a natural equivalence
\begin{equation}\label{v4-arith:md:eq:interfaces}
 L_H\otimes_{E_B}^{\mathbb L}L_K
 \xrightarrow{\sim}L_{K\otimes_B^{\mathbb L}H},
 \qquad g\otimes f\longmapsto g\circ(f\otimes1_H).
\end{equation}
It respects iterated composition. The identity $A$-bimodule
is sent to the regular $E_A$-bimodule.
\end{proposition}

\begin{proof}
The formula is balanced because acting on $g$ by precomposition
has the same effect as acting on $f$ by postcomposition.
Put $Q=P_A\otimes_A^{\mathbb L}K$. Both sides are exact
contravariant functors of $Q\in\operatorname{thick}_B(P_B)$.
At $Q=P_B$, the map is the unit isomorphism for tensor over
$E_B$.

It is therefore an equivalence on the thick closure.
For three interfaces, both parenthesizations compose the same
three maps. The derived associativity constraints give the same
comparison. For $K=A$, the definition gives
$L_A=\operatorname{RHom}_A(P_A,P_A)=E_A$ with its two regular
actions.
\end{proof}

The augmentation boundary and the identity interface have different
roles. The first becomes a regular dual module in
Theorem~\ref{v4-arith:md:thm:marked}. The second becomes the unit
bimodule for balanced composition. Replacing that bimodule by its
augmentation would change the composition law. A field realization
must preserve the chosen boundary action as well as these units.

\section{Comodules, contramodules and the direction of a map}
\label{v4-arith:md:sec:co-contra}

A coalgebra has a second correspondence that is covariant in its
module argument. Let $C$ be a counital coalgebra over a field $k$.
A left comodule has a coaction $M\to C\otimes_kM$.

A \emph{left contramodule} has instead a map
\[
 \pi:\operatorname{Hom}_k(C,P)\longrightarrow P.
\]
It satisfies the unit identity induced by the counit of $C$ and
the associativity identity obtained from
$\operatorname{Hom}_k(C,\operatorname{Hom}_k(C,P))$.
The two maps in that identity are composition with $\pi$ and
precomposition with the coproduct, using tensor--Hom adjunction.
For an infinite coalgebra this retains operations on all linear
maps from $C$, not only finite linear combinations.

For a right comodule $N$ and a left contramodule $P$, their
\emph{contratensor product} $N\odot_C P$ is the cokernel of
\[
 N\otimes_k\operatorname{Hom}_k(C,P)
 \rightrightarrows N\otimes_kP.
\]
One map is $n\otimes f\mapsto n\otimes\pi(f)$.
The other is $n\otimes f\mapsto\sum n_{(0)}\otimes f(n_{(1)})$,
where $\delta n=\sum n_{(0)}\otimes n_{(1)}$.
The sums here are finite on each comodule element. The subscript
$C$ distinguishes this quotient from the transported tensor
$\odot$ of \eqref{v4-arith:md:eq:dual-tensor}.

The regular coalgebra is a bicomodule. It defines functors
\begin{equation}\label{v4-arith:md:eq:co-contra}
 \Phi_C(P)=C\odot_C P,
 \qquad \Psi_C(M)=\operatorname{Hom}_C(C,M),
 \qquad \Phi_C\dashv\Psi_C.
\end{equation}
The Hom consists of colinear maps.

The left coaction on
$\Phi_C(P)$ comes from the first factor of the coproduct.
The right coaction on $C$ induces the contramodule operation
on $\Psi_C(M)$. In ungraded notation it sends
$F:C\to\operatorname{Hom}_C(C,M)$ to
$c\mapsto\sum F(c_{(2)})(c_{(1)})$.
Tensor--Hom adjunction identifies the two balancing conditions
and proves the adjunction in \eqref{v4-arith:md:eq:co-contra}.

For a dg coalgebra, use the same constructions in complexes with
Koszul signs. In the homotopy category of dg comodules, take the
thick subcategory generated by total complexes of exact triples
and closed under direct sums. The quotient is the
\emph{coderived category}.

The objects in this subcategory are called coacyclic.
The \emph{contraderived category} uses exact triples of
contramodules and closure under products. Its corresponding
objects are contraacyclic. These are the acyclic classes relevant
to the unbounded correspondence.

\begin{theorem}[The field-coefficient correspondence]
\label{v4-arith:md:thm:co-contra}
For a dg coalgebra over a field, the derived functors of
\eqref{v4-arith:md:eq:co-contra} give inverse covariant equivalences
\[
 \mathbb L\Phi_C:D^{\mathrm{ctr}}(C\text{-contra})
 \rightleftarrows D^{\mathrm{co}}(C\text{-comod}):\mathbb R\Psi_C.
\]
\end{theorem}

\begin{proof}
The coderived category is modeled by dg comodules that are
injective after forgetting the differential. The contraderived
category is modeled by contramodules that are projective after
forgetting the differential. These resolution statements are
Positselski, \emph{Two kinds of derived categories}, Theorem 4.4.
They concern the graded comodule and contramodule categories.

For a graded vector space $V$, the functors send the free
contramodule $\operatorname{Hom}_k(C,V)$ to the cofree comodule
$C\otimes_kV$, and conversely. The comparisons follow from
evaluation, coproduct and counit. Graded projectives are
summands of free contramodules, and graded injectives are
summands of cofree comodules. Thus the comparisons extend to
these two classes, commute with their differentials and give
inverse functors on the two homotopy models. This is the
correspondence theorem of that source, Section 5.2.
\end{proof}

For $C=k$, the functor $\Psi_C$ is the identity on complexes.
It is not vector-space dualization. Even for a finite coalgebra,
colinear Hom cannot be replaced by all linear maps.

Let
$C=k1\oplus kc$, with $c$ primitive. A colinear endomorphism
has the form
\[
 f(1)=a1,\qquad f(c)=b1+ac.
\]
Indeed the coaction identity first forces the displayed value
on $1$ and then the coefficient of $c$. Thus
$\dim_k\operatorname{Hom}_C(C,C)=2$, whereas
$\dim_k\operatorname{Hom}_k(C,C)=4$.

The injective condition also retains the coaction. The inclusion
$k1\subset C$ has no comodule retraction. If $r(1)=1$, the
colinearity equation on $c$ would read
$1\otimes r(c)=1\otimes r(c)+c\otimes1$. Thus the trivial
comodule $k$ is not injective, although its underlying vector
space is injective. Underlying coefficient injectivity alone
cannot justify the resolutions in the theorem.

\subsection{Bar and cobar supply additional algebra comparisons}

Let $A$ be an augmented dg algebra over $k$, and let $BA$ be
its normalized bar coalgebra. In this subsection use left
$A$-modules. Write $\tau:BA\to A$ for the
canonical degree-one twisting cochain.

Its identity is
\[
 d_A\tau+\tau d_{BA}+\mu_A(\tau\otimes\tau)\Delta=0,
\]
where the tensor of maps includes its Koszul sign.
For a module $M$, the twisted tensor $BA\otimes_\tau M$
has the ordinary tensor differential plus the term applying
the coproduct, $\tau$ and the module action. The identity above
makes its differential square to zero. Twisted Hom is defined
by the adjoint formula.

The resulting equivalences are
\begin{equation}\label{v4-arith:md:eq:bar-categories}
 D(A\text{-mod})\simeq D^{\mathrm{co}}(BA\text{-comod}),
 \qquad D(A\text{-mod})\simeq D^{\mathrm{ctr}}(BA\text{-contra}).
\end{equation}
They are Positselski's Theorem 6.3. One composite is the reduced
bar resolution of $M$. Its augmentation is a quasi-isomorphism.

For the other composite, filtration by the number of bar
entries reduces the unit comparison to the cofree coalgebra
resolution. The corresponding Hom argument uses products and
the contraderived acyclic class. Section 6.4 gives the analogous
statement for a conilpotent dg coalgebra $C$ and $A=\Omega C$.
The normalized algebra bars and their signs were constructed in
Chapter~\ref{v4-ch:bar-cobar-topology}.

The two equivalences in \eqref{v4-arith:md:eq:bar-categories}
intertwine the comodule--contramodule correspondence. On a
twisted cofree object this is the explicit identity
\[
 \Psi_C(C\otimes_\tau M)
 \cong\operatorname{Hom}_\tau(C,M).
\]
Thus coefficient transposition, bar--cobar comparison and
comodule--contramodule correspondence have specified connecting
maps. They remain different functors. Replacing the field by
$\mathbb Z_p[[T]]$ requires an exact structure and resolution
hypotheses over that ring.

\subsection{Changing the coalgebra}

Let $f:C\to D$ be a coalgebra map. Corestriction sends a left
$C$-comodule to a left $D$-comodule by
\begin{equation}\label{v4-arith:md:eq:corestriction}
 M\xrightarrow{\delta_M}C\otimes M
 \xrightarrow{f\otimes1}D\otimes M.
\end{equation}
It has the same underlying vector space. A $C$-contramodule
likewise becomes a $D$-contramodule by
\[
 \operatorname{Hom}_k(D,P)\xrightarrow{\operatorname{Hom}(f,P)}
 \operatorname{Hom}_k(C,P)\xrightarrow{\pi}P.
\]
Both functors go from $C$ to $D$. Applying the word
``opposite'' to $f$ does not produce a linear map $D\to C$.

The reverse comodule functor uses cotensor. Give $C$ its
$(C,D)$-bicomodule structure through $f$. For a left
$D$-comodule $N$, set
\[
 f^!N=C\mathbin{\square}_DN
 =\ker\bigl(C\otimes N\rightrightarrows C\otimes D\otimes N\bigr).
\]
The maps use the right $D$-coaction on $C$ and the left
$D$-coaction on $N$.

It is right adjoint to corestriction.
For example, the adjoint of $u:M\to N$ is
$m\mapsto\sum m_{(-1)}\otimes u(m_{(0)})$.
The inverse adjunction map applies the counit of $C$.
These formulas prove the adjunction identities.

On contramodules the reverse functor is
$f_!Q=\operatorname{Cohom}_D(C,Q)$. For a left $D$-comodule
$L$, define $\operatorname{Cohom}_D(L,Q)$ as the cokernel of
two maps
\[
 \operatorname{Hom}_k(D\otimes L,Q)
 \rightrightarrows\operatorname{Hom}_k(L,Q).
\]
One precomposes with the coaction of $L$. The other uses
\[
 \operatorname{Hom}_k(D\otimes L,Q)
 \cong\operatorname{Hom}_k(L,\operatorname{Hom}_k(D,Q))
\]
and then the contraaction of $Q$.
For $L=C$, the right $C$-coaction induces the left
$C$-contramodule structure on the quotient.
This is left adjoint to the contramodule corestriction.

\begin{proposition}[Naturality with derived adjoints]
\label{v4-arith:md:prop:naturality}
For dg coalgebras over a field there is a natural comparison
\begin{equation}\label{v4-arith:md:eq:naturality}
 \mathbb R\Psi_C\,\mathbb Rf^!
 \simeq\mathbb Lf_!\,\mathbb R\Psi_D
 :D^{\mathrm{co}}(D\text{-comod})
 \longrightarrow D^{\mathrm{ctr}}(C\text{-contra}).
\end{equation}
Here the derived functors use the graded injective and projective
models of Theorem~\ref{v4-arith:md:thm:co-contra}.
\end{proposition}

\begin{proof}
For a graded-injective $D$-comodule $N$, the comparison is
induced by
\[
 \operatorname{Hom}_C(C,C\mathbin{\square}_DN)
 \cong\operatorname{Hom}_D(C,N)
 \cong\operatorname{Cohom}_D(C,\operatorname{Hom}_D(D,N)).
\]
The first isomorphism is the cotensor adjunction. For a cofree
$N=D\otimes V$, the second is the evaluation identification
with $\operatorname{Hom}_k(C,V)$.

It extends to injective
summands. The cotensor functor preserves graded injectives,
and its cohomomorphism counterpart preserves graded
projectives. Thus the isomorphisms commute with differentials
and pass to the two homotopy models. This is the naturality
comparison of Positselski, Section 5.4.
\end{proof}

There is a related comparison when algebra maps and coalgebra
maps preserve a specified acyclic twisting cochain. If
$\tau:C\to A$, $\sigma:D\to B$, $g:A\to B$ and $f:C\to D$
satisfy $\sigma f=g\tau$, then extension over $g$ corresponds,
under the twisted tensor equivalences, to corestriction over
$f$. Restriction over $g$ corresponds to derived cotensor.
The chain identities follow by substituting $\sigma f=g\tau$
in the twisted differentials. See Positselski, Section 6.9.
The equation is an additional compatibility of actual maps.

\begin{example}[A linear splitting does not reverse a coaction]
\label{v4-arith:md:ex:retraction}
Let $D$ have basis $1,c,d$ and coproduct
\[
 \Delta c=1\otimes c+c\otimes1,
 \qquad\Delta d=1\otimes d+c\otimes c+d\otimes1.
\]
The counit kills $c,d$. Coassociativity follows by expanding
these formulas, or by dualizing multiplication in $k[x]/(x^3)$.
The subspace $C=k1\oplus kc$ is a subcoalgebra.

The linear retraction $r:D\to C$ fixing $1,c$ and killing $d$
preserves the counit, but
\[
 \Delta_Cr(d)=0,\qquad(r\otimes r)\Delta_Dd=c\otimes c\ne0.
\]
It is not a coalgebra map. Applying $1\otimes r$ to the right
regular coaction of $D$ gives a map $\delta:D\to D\otimes C$
whose coassociativity defect on $d$ is
\[
 (\delta\otimes1)\delta(d)-(1\otimes\Delta_C)\delta(d)
 =1\otimes c\otimes c.
\]
Thus even a counital splitting of underlying modules does not
give the reverse corestriction functor.
\end{example}

For a finite projective coalgebra map, the algebra transpose
$f^\vee:D^\vee\to C^\vee$ does exist. Ordinary restriction
of modules along it goes from $C^\vee$-modules to
$D^\vee$-modules. Extension in the other direction changes
the underlying module. The evaluation identity
\eqref{v4-arith:md:eq:transpose} remains valid without implying
any of the incorrect reverse functors.

\section{Completion and detection of a comparison}
\label{v4-arith:md:sec:completion}

Integral duality requires control of the entire coefficient tower.
Let $R$ be a commutative Noetherian ring and $I\subset R$ an
ideal. The derived $I$-adic completion of $M\in D(R)$ is
\begin{equation}\label{v4-arith:md:eq:derived-completion}
 \mathbb L\Lambda_I M
 =\operatorname{Rlim}_n(M\otimes_R^{\mathbb L}R/I^n).
\end{equation}
A complex is \emph{derived $I$-complete} when its natural map
to \eqref{v4-arith:md:eq:derived-completion} is an equivalence.
For Noetherian $R$ this is the reflection into the complexes
satisfying $\operatorname{RHom}_R(R[1/f],M)=0$ for every
$f\in I$. The equivalence of these definitions is the
Noetherian derived-completion theorem, Stacks Project, Tag 0922.

Write $K_n=M\otimes_R^{\mathbb L}R/I^n$.
The inverse-limit term is visible in the exact sequence
\begin{equation}\label{v4-arith:md:eq:milnor}
 0\longrightarrow\varprojlim_n{}^1H^{j-1}(K_n)
 \longrightarrow H^j(\mathbb L\Lambda_I M)
 \longrightarrow\varprojlim_nH^j(K_n)\longrightarrow0.
\end{equation}
Indeed the derived limit is the fiber of
$1-s:\prod_nK_n\to\prod_nK_n$, where $s$ applies the
transition to the next coordinate. Products of modules are
exact. The long exact cohomology sequence therefore identifies
the two outer terms with the cokernel and kernel of $1-s$.

The first derived limit vanishes for a surjective cohomology
tower, since $(1-s)x=y$ can then be solved recursively.
It also vanishes under the Mittag--Leffler condition: for each
fixed index the images of all sufficiently late transition
maps agree. This is weaker than surjectivity of each transition.

The reduction tower of
$M=[\mathbb Z_p\xrightarrow{p}\mathbb Z_p]$, in degrees zero
and one, gives an example:
\[
 H^0(M\otimes_{\mathbb Z_p}^{\mathbb L}\mathbb Z/p^n)
 =p^{n-1}\mathbb Z_p/p^n\cong\mathbb F_p.
\]
Every transition on this group is zero. Its eventual image is
zero, so it is Mittag--Leffler but is not a surjective tower.
Each term of the complex is flat over $\mathbb Z_p$.

Ordinary quotients can lose all the cohomology in the answer.
For $D_p=\mathbb Q_p/\mathbb Z_p$, one has $D_p/p^nD_p=0$,
but the two-term resolution of $\mathbb Z/p^n$ gives
\begin{equation}\label{v4-arith:md:eq:divisible-completion}
 D_p\otimes_{\mathbb Z_p}^{\mathbb L}\mathbb Z/p^n
 \simeq D_p[p^n][1],\qquad
 \mathbb L\Lambda_{(p)}D_p\simeq\mathbb Z_p[1].
\end{equation}
The transition on the kernels is multiplication by $p$.
Writing a kernel element as $a/p^n$ identifies it with the
reduction $\mathbb Z/p^{n+1}\to\mathbb Z/p^n$.
This tower is surjective, with inverse limit $\mathbb Z_p$.
Our shift convention places the displayed cohomology in
degree $-1$.

For finite cohomology modules there is a useful simplification.
If every $H^j(M)$ is finite over $R$, then
\[
 H^j(\mathbb L\Lambda_I M)
 \cong\varprojlim_nH^j(M)/I^nH^j(M).
\]
This is Stacks Project, Tag 0A06.

The Noetherian
Artin--Rees theorem makes completion exact on finite modules.
Finite free resolutions and truncations give the assertion
for complexes. In particular finite cohomology over
$\mathbb Z_p[[T]]$ is already $p$-adically complete.
This does not justify replacing the derived tensor in
\eqref{v4-arith:md:eq:derived-completion} by every chosen
termwise quotient.

\subsection{The unit and counit must be the maps tested}

\begin{proposition}[Conservativity of reduction]
\label{v4-arith:md:prop:conservative}
Let $u:M\to N$ be a map of derived $I$-complete complexes.
If $u\otimes_R^{\mathbb L}R/I$ is an equivalence, then $u$
is an equivalence.
\end{proposition}

\begin{proof}
The cone $K$ is derived complete, since derived completion is
exact. The hypothesis says $K\otimes_R^{\mathbb L}R/I=0$.
It follows that $K\otimes_R^{\mathbb L}V=0$ for every
$R/I$-module $V$, by associativity of derived tensor.
The finite filtration of $R/I^n$ by powers of $I$ has such
successive quotients. Induction therefore gives
$K\otimes_R^{\mathbb L}R/I^n=0$ for every $n$.
Taking the derived limit and using completeness gives $K=0$.
\end{proof}

Suppose now that an integral adjunction $F\dashv G$ is defined
between stable coefficient categories. Suppose both categories
have exact reduction functors, conservative on the complete
objects under consideration. Assume $F,G$ preserve these
objects and that their reductions form an adjunction
$\bar F\dashv\bar G$. The comparison must identify the
reductions of the actual unit and counit with the unit and
counit of $\bar F\dashv\bar G$.

Under these hypotheses, if the reduced adjunction is an
equivalence, then the original adjunction is an equivalence
on that complete domain. Its unit and counit have zero
reduced cones, so conservativity makes both maps equivalences.
For coderived or contraderived recipients, this argument
requires an internal detection statement for those cones.
An underlying ordinary derived complex suffices only when the
forgetful comparison detects them.

A compatible tower gives another positive construction.
Let $\mathcal C_n\simeq\mathcal D_n$ be equivalences of
categories, equipped with coherent comparisons to their
transition functors. They form an equivalence of diagrams,
so induce
\[
 \varprojlim_n\mathcal C_n\simeq\varprojlim_n\mathcal D_n.
\]
The limits retain their descent data and mapping spaces.
Equivalences chosen independently at each level do not yet
give the required equivalence of diagrams.

\subsection{Limits of reduction}

Conservativity detects equivalences of supplied maps. It is
not faithfulness: multiplication by $p$ on $\mathbb Z_p$ is
nonzero and reduces to zero. Nor does reduction make a lift
unique. For $R=\mathbb Z_p[[T]]$, the modules
\[
 M_0=R/(T),\qquad M_1=R/(T-p)
\]
have the same reduction modulo $p$. They are not isomorphic
as $R$-modules, since their annihilator ideals differ.
Even a fixed lift has distinct automorphisms $1$ and $1+p$
with the same reduction.

There is a second distinction between scalar and power-series
coefficients. The module $R/(T)$ is free over $\mathbb Z_p$,
but the free $R$-resolution
$[R\xrightarrow{T}R]$ gives
\[
 \operatorname{Tor}_1^R(R/(T),R/(T))=R/(T).
\]
It is not flat over $R$. Its Iwasawa invariant $\mu$ is zero,
because $T$ becomes invertible at the height-one prime $(p)$.
Thus $\mu=0$ supplies neither flatness nor the required
base-change comparison.

The coefficient ring $\Lambda=\mathbb Z_p[[T]]$ is regular
local of dimension two. The parameters $p,T$ form a regular
sequence, and their Koszul complex resolves the residue field
in length two. A minimal free resolution of a finite module has
zero differentials after reduction to the residue field. Its
terms then compute Tor. Computing Tor instead with the Koszul
resolution shows that every term above degree two is zero,
by Nakayama's lemma.

Finite $\Lambda$-modules therefore have
projective dimension at most two. In particular this holds
for $\Lambda/J$ for every ideal $J$, since $\Lambda$ is
Noetherian. Dimension shifting and Baer's injectivity criterion
then bound the injective dimension of every $\Lambda$-module
by two.
These bounds can support resolution arguments, but the exact comodule
and contramodule structures must still be specified.
In particular the module condition $\mu=0$ and coalgebra
conilpotence are independent, as the two explicit coalgebras
in Proposition~\ref{v4-arith:coefficients:prop:mu} show.

For an abelian number field, form the Iwasawa module from
the $p$-primary class groups of its cyclotomic tower.
Ferrero--Washington proves that its $\mu$ invariant vanishes.
Transferring this invariant to another
arithmetic object requires a comparison of the specified
coefficient modules, for example a pseudo-isomorphism.
Neither the oscillator basis nor the coproduct alone provides
that comparison.

\section{Completed coordinate coactions and Fock fields}
\label{v4-arith:md:sec:completed-modules}

The transpose of a coaction can have infinitely many output
coordinates even when each input state is a polynomial. This
does not prevent a continuous coordinate coaction. It does
force a choice of recipient.

Let $k$ be a discrete field of characteristic zero and let
$H$ be a Hopf algebra with bijective antipode $\mathsf s$.
Suppose a right comodule is an increasing union of finite
subcomodules, $M=\bigcup_NM_N$. Define
\begin{equation}\label{v4-arith:md:eq:coordinate-limits}
 D_M=\varprojlim_NM_N^*,\qquad
 \mathscr T_r(M)=\varprojlim_N(M_N^*\otimes_kH^{\otimes r}).
\end{equation}

The finite-stage spaces are discrete. An element of
$\mathscr T_r(M)$ is equivalently a linear map
$M\to H^{\otimes r}$, with the topology of restriction to
finite stages. In particular the target of a completed
coaction is $\mathscr T_1(M)$. It need not be
$D_M\otimes_kH$.

Use the following tensor convention on right comodules:
\begin{equation}\label{v4-arith:md:eq:reversed-tails}
 \delta(m\otimes n)
 =m_{(0)}\otimes n_{(0)}\otimes n_{(1)}m_{(1)}.
\end{equation}
The order of the two Hopf-algebra coefficients is part of the
convention. Associativity of $H$ makes this a monoidal tensor.

\begin{proposition}[The continuous coordinate coaction]
\label{v4-arith:md:prop:completed-coaction}
There is a natural continuous map
\begin{equation}\label{v4-arith:md:eq:completed-coaction}
 \widehat\delta_M:D_M\longrightarrow\mathscr T_1(M),
 \qquad (\widehat\delta_Mf)(m)
 =\sum f(m_{(0)})\mathsf s^{-1}(m_{(1)}).
\end{equation}
It is counital and coassociative in the corresponding
coordinate recipients. It gives a natural tensor comparison
\[
 D_{M\otimes N}\cong
 D_N\widehat\otimes D_M
 :=\varprojlim_{a,b}(N_b^*\otimes M_a^*),
\]
where both sides retain their inverse systems of comodules.
\end{proposition}

\begin{proof}
In a basis $e_j$ of $M_N$, write
$\delta e_j=\sum_i e_i\otimes u_{ij}$. The coaction identity
is $\Delta u_{ij}=\sum_lu_{il}\otimes u_{lj}$.
The proposed dual coaction sends
$e^i$ to $\sum_je^j\otimes\mathsf s^{-1}(u_{ij})$.
Since the antipode reverses coproduct,
\[
 \Delta\mathsf s^{-1}(u_{ij})
 =\sum_l\mathsf s^{-1}(u_{lj})\otimes\mathsf s^{-1}(u_{il}).
\]

This is precisely coassociativity of the dual formula.
The counit identity follows from the identity coefficient
matrix. Restriction to a smaller subcomodule commutes with
the formula, so it defines the inverse-limit map. Each
finite projection uses only $f|_{M_N}$, which proves continuity.

At a finite pair of stages, $g\otimes f$ corresponds to the
functional $m\otimes n\mapsto f(m)g(n)$. Its Hopf coefficient
is compatible because
\[
 \mathsf s^{-1}(n_{(1)}m_{(1)})
 =\mathsf s^{-1}(m_{(1)})\mathsf s^{-1}(n_{(1)}).
\]
This is the coefficient of $g\otimes f$ under
\eqref{v4-arith:md:eq:reversed-tails}. Passing through the
restriction maps gives the tensor comparison. A map of the
exhaustive filtered unions takes each finite-dimensional
stage into a finite stage, so its transpose is continuous
and all comparisons are natural.
\end{proof}

The reversal is visible in the four-dimensional Hopf algebra
generated by $g,x$, with
\[
 g^2=1,\quad x^2=0,\quad gx=-xg,
 \qquad\Delta g=g\otimes g,\quad
 \Delta x=1\otimes x+x\otimes g.
\]
Its inverse antipode has $\mathsf s^{-1}(x)=-gx$.
Consequently $\mathsf s^{-1}(xg)=-x$, while
$\mathsf s^{-1}(x)\mathsf s^{-1}(g)=x$.
Keeping the original tensor order changes this coefficient
by $2x$.

\subsection{The full binomial coaction}

Take $F_K=k[x_1,x_2,\ldots]$, with $K\ne0$, and give $x_n$
energy $n$. Let $H=k[b_2,b_3,\ldots]$ with every $b_j$
primitive, and set
\[
 \delta_K(x_n)=x_n\otimes1-nK\,1\otimes b_{n+1}.
\]
Extending as an algebra map gives a right coaction. The
subspace of energy at most $N$ is finite-dimensional and
stable under it.

Let $f_\alpha$ extract the coefficient of
$x^\alpha=\prod_jx_j^{\alpha_j}$, where $\alpha$ has finite
support. Applying \eqref{v4-arith:md:eq:completed-coaction}
to the binomial expansion gives
\begin{equation}\label{v4-arith:md:eq:fock-coaction}
 \widehat\delta_K(f_\alpha)
 =\sum_{\beta\geq\alpha}\binom\beta\alpha f_\beta\otimes
       \prod_{j\geq1}(jKb_{j+1})^{\beta_j-\alpha_j},
 \qquad \binom\beta\alpha=\prod_j\binom{\beta_j}{\alpha_j}.
\end{equation}
Each finite-energy projection has finitely many terms.
For the vacuum functional,
$(\widehat\delta_Kf_0)(x_1^n)=(Kb_2)^n$.
The arbitrarily high powers show why the algebraic
restricted-dual tensor is too small for this formula.

The construction also works over a formal coefficient tower.
Suppose that $H_d$ and $M_{N,d}$ are compatible Hopf algebras
and finite free subcomodules over $k[\hbar]/(\hbar^d)$,
with bijective antipodes. The same finite matrix proof gives
a coaction with target
\[
 \varprojlim_{N,d}
 \bigl(M_{N,d}^{\vee}\otimes_{k[\hbar]/(\hbar^d)}H_d\bigr).
\]
For example, a translation derivation satisfying
$Tb_n=nb_{n+1}$ gives primitive tails
\[
 \beta_{K,n}=\frac{b_n-e^{K\hbar T}b_n}{\hbar}
 =-\sum_{r\geq1}\frac{K^r\hbar^{r-1}}{r!}T^rb_n.
\]

This quotient is defined by the series and does not invert
$\hbar$. On the polynomial Hopf algebra with primitive
$b_n$, these tails define
$\delta_K(x_n)=x_n\otimes1+1\otimes\beta_{K,n}$.
At each precision only finitely many terms of each tail
remain. The limit recovers a continuous coordinate coaction
on $\prod_\alpha k[[\hbar]]f_\alpha$.

\subsection{Evaluation imposes another topology}

Let $M=\bigoplus_{n\geq0}M[n]$, with each $M[n]$ finite,
and let $D^{\mathrm{res}}=\bigoplus_nM[n]^*$.
Consider the matrix completion
\begin{equation}\label{v4-arith:md:eq:matrix}
 \mathscr E(M)=\prod_{p,q\geq0}M[p]\otimes M[q]^*,
\end{equation}
with discrete finite blocks. The finite sums of the identity
tensors on $M[0],\ldots,M[N]$ converge to an element of
$\mathscr E(M)$ with identity diagonal blocks and zero
off-diagonal blocks. This is the formal diagonal coevaluation.

\begin{proposition}[Failure of continuous contraction]
\label{v4-arith:md:prop:evaluation}
If infinitely many $M[n]$ are nonzero, ordinary evaluation
on $M\otimes D^{\mathrm{res}}$ has no continuous extension
to \eqref{v4-arith:md:eq:matrix}.
\end{proposition}

\begin{proof}
Choose distinct $n_j\to\infty$, vectors $e_j\in M[n_j]$
and functionals $f_j$ with $f_j(e_j)=1$.
The tensors $e_j\otimes f_j$ converge to zero in the
product topology, since every fixed finite collection of
blocks is eventually zero. Their evaluations are all $1$,
which does not converge to zero in $k$. This contradicts
continuity.
\end{proof}

The argument persists over $k[[\hbar]]$ with the product of
the $\hbar$-adic block topologies: the constant scalar $1$
still does not approach zero. Exchanging the two tensor
orders is a homeomorphism of the underlying matrix products,
so it does not repair continuity. No assertion about that
flip as a comodule map is needed.

Weighting the coevaluation by $q^n$ gives a different map.
Its zigzag on $M[n]$ is $q^n\operatorname{id}$ rather than
$\operatorname{id}$. Such a weighted sewing operator can
be useful, but it is not the coevaluation of ordinary
duality. The coordinate coaction above and the diagonal
coevaluation therefore do not make this matrix completion
into a rigid tensor category.

\subsection{The inverse coordinate and the central level}

On $F_K$ use the Heisenberg modes
$J_{-n}=x_n$, $J_n=nK\partial_{x_n}$ for $n>0$, and
$J_0=0$. Then $[J_m,J_n]=mK\delta_{m+n,0}$.
For an operator $A$, its transpose is
$(A^tf)(v)=f(Av)$.
At the same coordinate the evaluation-compatible Lie
action is $-J_m^t$. It has level $-K$. It is not lower
truncated on the restricted dual: the positive modes
send $f_0$ to $-nKf_{e_n}$ for every $n>0$.

In the inverse coordinate $w=z^{-1}$, the weight-one
contragredient instead has
\begin{equation}\label{v4-arith:md:eq:inverse-coordinate}
 J^\infty(w)=-w^{-2}J(w^{-1})^t,
 \qquad J^\infty_m=-J_{-m}^t.
\end{equation}
This is the primary-current specialization of Huang,
\emph{Introduction to vertex operator algebras, III},
equation (1.26). Its properties here can be checked directly.

\begin{proposition}[Restricted contragredient Fock module]
\label{v4-arith:md:prop:contragredient}
The current \eqref{v4-arith:md:eq:inverse-coordinate} is
lower truncated on $F_K^{\mathrm{res}}=\bigoplus_nF_K[n]^*$.
It has level $K$ and is isomorphic to $F_K$ by
\begin{equation}\label{v4-arith:md:eq:fock-identification}
 x^\alpha\longmapsto
 \left(\prod_j(-jK)^{\alpha_j}\alpha_j!\right)f_\alpha.
\end{equation}
On the full product $\prod_nF_K[n]^*$ the same current
need not be lower truncated.
\end{proposition}

\begin{proof}
For $n>0$ its modes are
\[
 J^\infty_nf_\alpha=-f_{\alpha-e_n},\qquad
 J^\infty_{-n}f_\alpha=-(\alpha_n+1)nKf_{\alpha+e_n}.
\]
The first expression is zero when $\alpha_n=0$.
Positive modes lower dual energy, so sufficiently large
ones kill every vector in the direct sum. These formulas
also prove \eqref{v4-arith:md:eq:fock-identification}, whose
coefficients are nonzero because $k$ has characteristic
zero and $K\ne0$.
Transposition reverses commutators, giving
\[
 [J^\infty_m,J^\infty_n]
 =-[J_{-m},J_{-n}]^t=mK\delta_{m+n,0}.
\]
Normal ordering and differentiation then identify all the
Heisenberg module fields.

For the full product, take the functional $f$ with
$f(x_n)=1$ for every $n>0$ and zero on the other monomials.
Its vacuum component is
$(J^\infty(w)f)(1)=-\sum_{n\geq1}w^{-n-1}$.
It has no lower Laurent bound, proving the last assertion.
\end{proof}

The inverse coordinate restores truncation on the restricted
dual and restores the original level. For a tensor that adds
levels and has scalar unit of level zero, it is therefore
not an inverse-level dual when $K\ne0$.
The field-module isomorphism does not assert compatibility
with the completed coordinate coaction. Finding a category
in which the desired evaluation, coaction and field maps
coexist remains a further construction. The two explicit
failures above concern their stated recipients.

\section{Full collision coefficients}
\label{v4-arith:md:sec:full}

A singular field needs a coefficient category that retains
all its translates. Suppose $Y(a,z)b$ has a nonzero pole
over a field of characteristic zero. A relation $p(T)a=0$
would give $p(\partial_z)Y(a,z)b=0$. If the lowest Laurent
term is $vz^{-q}$ and $p$ has leading term $ct^d$, the
coefficient of $z^{-q-d}$ is
\[
 c(-q)(-q-1)\cdots(-q-d+1)v\ne0.
\]
No lower derivative contributes there. Thus the translates
$T^na$ are linearly independent.

On $X=\mathbb A^1_k$, the right differential-module action is
\begin{equation}\label{v4-arith:md:eq:right-current}
 (f\,dz\otimes a)\partial_z
 =-f'\,dz\otimes a-f\,dz\otimes Ta.
\end{equation}
The translation orbit therefore gives a free rank-one
$\mathcal D_X$-submodule of $\omega_X\otimes V$.
It is perfect over $\mathcal D_X$, although it is not
holonomic: its characteristic variety is the full cotangent
bundle. Holonomicity, perfectness and dualizability for a
collision tensor are different conditions.

At finite formal precision a nonzero pole can have a torsion
coefficient. Over $k[\hbar]/(\hbar^2)$ the state $\hbar J$
can have a nonzero pole and is still killed by $\hbar$.
For a scalar-extended family, freeness of the translation
orbit over the parameter ring follows if the pole remains
nonzero modulo $\hbar$. A least-order coefficient of any
proposed relation then contradicts the argument over $k$.

\subsection{The retained perfect diagram}

Let $U$ be a smooth affine chart of dimension $d$, and set
$\mathcal D=\Gamma(U,\mathcal D_U)$. A perfect right
$\mathcal D$-complex is a summand of a finite complex of
finite free right modules. A coordinate volume form converts
left modules to right modules by
$f^t=f$, $\partial_i^t=-\partial_i$ and $(ab)^t=b^ta^t$.
Define
\begin{equation}\label{v4-arith:md:eq:perfect-differential-dual}
 P^\dagger=
 \bigl(\operatorname{RHom}_{\mathcal D}(P,\mathcal D)\bigr)^t[d].
\end{equation}

The canonical line gives the same conversion without a
chosen volume form. Finite projective evaluation proves
$(P^\dagger)^\dagger\simeq P$. The two shifts cancel because
duality reverses shifts. The same finite calculation commutes
with restriction to an open chart and with external products.
Dimensions add in the latter comparison.

Every derived $\mathcal D$-module has a filtered presentation
$M\simeq\varinjlim_\alpha P_\alpha$ by perfect objects.
The presentation is in the derived category, and need not
be a union of submodules. Define its full coefficient dual
and the inverse construction by
\begin{equation}\label{v4-arith:md:eq:full-dual}
 M^\sharp=\{P_\alpha^\dagger\}_\alpha,
 \qquad\mathbb S\{Q_\beta\}=\varinjlim_\beta Q_\beta^\dagger.
\end{equation}
The braces retain the pro-object. The fact used here is that
the derived module category is the ind-completion of its
perfect objects, since the regular module is a compact
generator.

\begin{proposition}[Full coefficient polarity]
\label{v4-arith:md:prop:full}
The functors in \eqref{v4-arith:md:eq:full-dual} give a
presentation-independent anti-equivalence
\[
 D(\mathcal D\text{-right})^{\mathrm{op}}
 \simeq\operatorname{Pro}(\operatorname{Perf}_{\mathrm r}(\mathcal D)).
\]
\end{proposition}

\begin{proof}
For $M=\varinjlim_\alpha P_\alpha$ and
$N=\varinjlim_\beta Q_\beta$, the mapping formula is
\begin{align*}
 \operatorname{Map}_{\mathrm{Pro}}(N^\sharp,M^\sharp)
 &\simeq\varprojlim_\alpha\varinjlim_\beta
       \operatorname{Map}(Q_\beta^\dagger,P_\alpha^\dagger)\\
 &\simeq\varprojlim_\alpha\varinjlim_\beta
       \operatorname{Map}(P_\alpha,Q_\beta)
 \simeq\operatorname{Map}(M,N).
\end{align*}
The last equivalence uses compactness of $P_\alpha$.
The formula respects composition. Applying it to any
pro-perfect diagram identifies that diagram with the dual
of its image under $\mathbb S$. It proves both inverse
comparisons and independence of the presentation.
\end{proof}

The pro-object contains more information than its ordinary
inverse limit. For the Weyl algebra
$D=k\langle u,\partial\rangle/([\partial,u]-1)$, the
right-module localization has presentation
\[
 D[u^{-1}]=\varinjlim
 (D\xrightarrow{L_u}D\xrightarrow{L_u}\cdots),
 \qquad a\text{ at stage }n\longmapsto u^{-n}a.
\]
Its coefficient dual is the tower of $D[1]$ with transitions
$L_u$. No transition to a fixed stage is zero, so this
pro-object is nonzero. Its underived inverse limit is zero:
each coordinate of a compatible sequence lies in
$\bigcap_nu^nD=0$, as follows from polynomial coefficients
in Weyl normal form. No vanishing of a derived inverse
limit is asserted.

The vacuum has another finite resolution,
$[D\xrightarrow{L_\partial}D]\to\omega_X$.
It gives $\omega_X^\dagger\simeq\omega_X$, whereas
$D^\dagger=D[1]$. Thus the full dual retains both the
vacuum and the free translation orbit with their respective
shifts.

\subsection{Transporting the collision tensor}

Use the ordered collision schemes $\mathscr R_I$ and their
exceptional collision pullbacks. A coefficient system
assigns a right differential complex to each scheme,
together with coherent comparisons along these pullbacks.
Let $\mathcal C$ denote this category. Define
$\mathcal C^\sharp$ by applying
Proposition~\ref{v4-arith:md:prop:full} on affine charts and
gluing. Opposite categories preserve limits, so the affine
comparisons give $\mathcal C^{\mathrm{op}}\simeq
\mathcal C^\sharp$ on the descent categories.

For any geometric functor $F$, its transported functor is
$F^\sharp=\sharp F\mathbb S$. In particular collision
pullback becomes $\delta^\dagger=\sharp\delta^!\mathbb S$,
and open direct image becomes
$j_{!,\mathrm{pro}}=\sharp j_*\mathbb S$.
The latter is a functor to pro-perfect coefficients. It is
not an assertion that ordinary $j_!$ exists on every large
differential-module coefficient in this form.

For an ordered decomposition $I=G\sqcup H$, let $U_{G|H}$
be the open where the two groups of points are disjoint,
and let $j_{G|H}$ be its inclusion. Internal collisions
within a group remain allowed. The atlas convolution
$\star$ sums the corresponding open direct images over
all ordered decompositions. Its transported tensor is
\begin{equation}\label{v4-arith:md:eq:full-tensor}
 (E\star_\sharp F)_I
 =\bigoplus_{I=G\sqcup H}j_{G|H,!,\mathrm{pro}}
 \left.(E_G\boxtimes^\sharp F_H)\right|_{U_{G|H}}.
\end{equation}

Here the external product uses the two-index perfect
diagrams. The sum is finite. Equivalently,
$E\star_\sharp F=(\mathbb SE\star\mathbb SF)^\sharp$.
Associativity and every collision comparison are transported
through the same coefficient duality.

The tensor unit $K$ has coefficient $k$ at empty support
and zero at nonempty support. Thus $K^\sharp=K$.
The word-cut tensor, whose decompositions are consecutive
prefixes and suffixes, is a separate construction.
In the atlas tensor, $13\mid2$ and $2\mid13$ are present.
Transposition retains these components and all normal
derivatives rather than replacing a full collision by an
ordered word.

For finite lists there is consequently an equivalence of
operation spaces
\begin{equation}\label{v4-arith:md:eq:operation-spaces}
 \operatorname{Map}_{\mathcal C}
       \left(\mathop{\star}_iX_i,\mathop{\star}_jY_j\right)
 \simeq
 \operatorname{Map}_{\mathcal C^\sharp}
       \left(\mathop{\star_\sharp}_jY_j^\sharp,
                         \mathop{\star_\sharp}_iX_i^\sharp\right).
\end{equation}
It follows directly from the mapping equivalence and
\eqref{v4-arith:md:eq:full-tensor}. It reverses composition
and retains identities, homotopies and their coherences.

\subsection{Associative algebras in the two polarities}

An augmented coefficient algebra is called connected here
when its empty coefficient is $k$. Its augmentation ideal
then has zero empty coefficient. On a configuration $I$,
a tensor word in that ideal has at most $|I|$ nonempty
factors. This support bound is the finiteness condition for
normalized bar and cobar, independently of the dimension
of any individual coefficient.

\begin{theorem}[Associative coefficient duality]
\label{v4-arith:md:thm:coefficient-koszul}
On the connected augmented domains for the full atlas
tensor, normalized bar and full coefficient duality give
an anti-equivalence
\[
 \operatorname{Alg}_{E_1}^{\mathrm{aug,conn}}
       (\mathcal C,\star)^{\mathrm{op}}
 \simeq
 \operatorname{Alg}_{E_1}^{\mathrm{aug,conn}}
       (\mathcal C^\sharp,\star_\sharp).
\]
Its functor is $\mathcal K(A)=(B_\star A)^\sharp$, and its
inverse is $E\mapsto\mathbb S B_{\star_\sharp}E$.
\end{theorem}

\begin{proof}
The bar coproduct transposes to the multiplication of
$\mathcal K(A)$ through \eqref{v4-arith:md:eq:full-tensor}.
The bar counit becomes its unit, and its coaugmentation
becomes its augmentation.
For a connected coalgebra $C$, dualizing the normalized
bar of $C^\sharp$ gives
$\mathbb S B_{\star_\sharp}(C^\sharp)\simeq\Omega_\star C$.
Suspensions reverse, multiplication transposes to the
coproduct, and outer deconcatenation becomes concatenation.

The bar--cobar comparison is finite on every $I$.
Filter a composite by its original positive-support entries.
For more than one entry, each gap is either an outer cut or
an inner product gap. On the associated graded these two
choices form a two-term identity complex.

Contract the
first gap, with the tensor sign. The single-entry comparison
is the identity, and the finite filtration proves the unit
and counit equivalences. This is the finite-support
bar--cobar argument of Chapter~\ref{v4-ch:bar-cobar-topology}
applied to all ordered decompositions.
Its triangles, together with coefficient biduality, give
the two inverse functors.
\end{proof}

For a formal family the same proof applies at every
$R_m=k[\hbar]/(\hbar^m)$ using relative differential
operators $\mathcal D_U\otimes_kR_m$.
Finite perfect duality does not require $R_m$ to be regular.
Keep the coefficient-reduction functors and their transported
pro-diagrams. No exchange of a pro-limit, an adic limit and
a Laurent expansion is used.

This theorem constructs a second associative coefficient
algebra with an inverse comparison. The Lie--commutative
chiral correspondence concerns other algebra objects and
does not supply this associative comparison. A realization
of its output by fields must additionally identify the
field products, exchange, tensor and local vacuum. The
next construction supplies such a realization with explicit
markings.

\section{A field realization with a bar contact marking}
\label{v4-arith:md:sec:charged}

The field data must be strong enough to recover the input
algebra from its dual. For charged currents the relevant
data are a quadratic current core, its translation and
jet derivations, and the contact operations retained in
the bar. We work over $k=\mathbb C$ at the fixed precision
$R_2=k[\hbar]/(\hbar^2)$. The parameter has degree, charge
and weight zero. Position on the affine line is denoted
by $z$ and field separation by $u$.

\subsection{Charged jets and their fields}

Let $E$ be a bounded finite-dimensional complex in
nonpositive cohomological degrees. Give $e_n=e\otimes t^n$
charge one and weight $n+1$. On
$\operatorname{Sym}_k(E\otimes k[t])$ define even
derivations
\begin{equation}\label{v4-arith:md:eq:jets}
 Te_n=(n+1)e_{n+1},\qquad
 \delta_i e_n=e_{n-i+1}\quad(i\geq1),\qquad e_j=0\ (j<0).
\end{equation}
They kill the unit, commute with the differential and satisfy
$[\delta_i,\delta_j]=0$ and
$[\delta_i,T]=(i-1)\delta_{i-1}$, with $\delta_0=0$.
The derivation $\delta_1$ measures charge.

A \emph{charged jet algebra} is a quotient of this dg
algebra by an ideal homogeneous in charge and weight and
stable under the displayed derivations. Require its
charge-zero part to be $k$. Each weight piece is then
a bounded finite-dimensional complex. The quotient
presentation is not an extra marking.

The charge-one generators are recovered intrinsically.
On $A_1$ put $D=\delta_2$. Then $[D,T]=1$ and $D$ is
locally nilpotent. The formulas
\begin{equation}\label{v4-arith:md:eq:seed}
 \pi=\sum_{j\geq0}\frac{(-1)^j}{j!}T^jD^j,
 \qquad
 a=\sum_{n\geq0}\frac{T^n}{n!}\pi(D^na)
\end{equation}
give $A_1=E_A\otimes k[t]$, where $E_A=\ker(D|_{A_1})$.

Both sums are finite on every input. Commuting $D$ through
$T^j$ cancels adjacent terms and proves $D\pi=0$.
The binomial identity proves the second formula. Successive
powers of $D$ prove its uniqueness. Thus $E_A$ and the
kernel of $\operatorname{Sym}(E_A\otimes k[t])\to A$ are
recovered from $A$ and its jet actions.

After scalar extension to $R_2$, define
\begin{equation}\label{v4-arith:md:eq:charged-fields}
 \begin{aligned}
 C_1(u)&=\sum_{p,q\geq0}(-1)^{p+1}\binom{p+q}{p}
 u^{-p-q-1}\delta_{p+1}\otimes\delta_{q+1},\\
 E(u)&=1-\hbar C_1(u),\qquad
 Y_A(a,u)b=\mu_A(e^{uT}\otimes1)E(u)(a\otimes b),\\
 \mathcal S_A(u)&=1+2\hbar C_1(u).
 \end{aligned}
\end{equation}
On fixed algebraic inputs the contraction is finite.
Taylor translation contributes nonnegative powers, so the
field output still has a finite lower Laurent bound.
The symbol $\mathcal S_A$ denotes unpermuted exchange.

An input morphism is an $R_2$-linear unital dg algebra
map preserving charge and commuting with $T$ and every
$\delta_i$. In particular it preserves multiplication,
$f\mu_A=\mu_{A'}(f\otimes f)$.
These equations imply that it preserves the fields and
exchange in \eqref{v4-arith:md:eq:charged-fields}.
The weak equivalences are state quasi-isomorphisms.

The first interacting one-seed quotient is
\begin{equation}\label{v4-arith:md:eq:cubic}
 A_0=k[e_0,e_1,\ldots]/I^3=k\oplus V\oplus Q,
 \qquad I=(e_0,e_1,\ldots),\quad Q=\operatorname{Sym}^2V.
\end{equation}
Its differential is zero. The field of $e_0$ satisfies
\[
 Y_{A_0}(e_0,u)e_0=(1+\hbar/u)\sum_{n\geq0}e_ne_0u^n.
\]
Its $\hbar u^{-1}$ coefficient is $e_0^2\ne0$.
The quotient by $I^2$ would kill this interaction.

\subsection{The marked current core}

Put $\mathfrak g=\mathfrak{sl}_2(k)$, with
$[h,e]=2e$, $[h,f]=-2f$ and $[e,f]=h$.
An independent current core consists of a discrete
nonpositive complex $L$ with bounded finite-dimensional
nonnegative weight pieces and completed generators
$W=(L[1])^\vee$. Here the dual of a discrete space is its
full degreewise dual with product topology. The dual of
a completed space is its continuous degreewise dual.
Thus $W^\vee[-1]=L$.

The finite-weight generator space is
$W_{\mathrm{pol}}=\bigoplus_q(L_q[1])^*$.
Its polynomial core is $\operatorname{Sym}(W_{\mathrm{pol}})$,
and its completion is the product of the weight pieces
in each cohomological degree. Require a continuous
square-zero derivation whose values on generators have
linear and quadratic parts. If $|v_i|=a_i$,
$d_Lv_i=\sum_kD_i^kv_k$ and $[v_i,v_j]=\sum_kf_{ij}^kv_k$,
our sign convention is
\begin{equation}\label{v4-arith:md:eq:current-differential}
 dc^k=\sum_i(-1)^{a_i}D_i^kc^i
 -\frac12\sum_{i,j}(-1)^{a_i(1-a_j)}f_{ij}^kc^ic^j.
\end{equation}
The linear, quadratic and cubic components of $d^2=0$
give the differential, derivation and Jacobi identities
on $L$.

Mark a scalar-current map dual to a dg Lie map
$j:\mathfrak g\to L$, and require translation to kill
$j(\mathfrak g)$. For $x\cdot v=[j(x),v]$, require
\begin{equation}\label{v4-arith:md:eq:isotypic}
 \mathfrak g\otimes
 \operatorname{Hom}_{\mathfrak g}(\mathfrak g,L)
 \xrightarrow{\sim}L.
\end{equation}
Charge, translation and all jet actions are retained.
The recovered multiplicity space must satisfy the
charged jet conditions above. This is a condition on
the supplied core and its operations.

For a charged jet algebra $A$, the forward core is
constructed from
\[
 L_A=\mathfrak g\otimes A,
 \quad[x\otimes a,y\otimes b]=[x,y]\otimes ab,
 \quad d_{L_A}=1\otimes d_A,
 \quad j_A(x)=x\otimes1.
\]
Its translation and jets are $1\otimes T$ and
$1\otimes\delta_i$. The corresponding derivations on
the dual generators are their negative transposes,
denoted $T_B,D_i$.

Let $\mathcal C^{(n)}$ be charge $n$ of the polynomial
core, and set
\begin{equation}\label{v4-arith:md:eq:current-localization}
 \begin{aligned}
 \Lambda_n&=k[x_1,\ldots,x_n]^{S_n},
 &d_n&=\prod_{i<j}(x_i-x_j)^2,\\
 B^{\mathrm{mer}}&=\bigoplus_{n\geq0}
 \mathcal C^{(n)}\otimes_{\Lambda_n}\Lambda_n[d_n^{-1}].
 \end{aligned}
\end{equation}
Here $d_0=d_1=1$. The variables $x_i$ are constituent
dual-jet variables, not positions on the affine line.
Their power sums act by $p_m\mapsto-D_{m+1}$, including
$p_0=n$. This symmetric-ring action is determined by
the marked jets.

Multilinear polarization explains the action and the
localization. At charge $n$, distribute the $n$ labeled
constituents among current generators. Their jet indices
give monomials in $x_1,\ldots,x_n$.

Take symmetric
invariants with cohomological Koszul signs.
Multiplication uses the normalized projector
$(n!)^{-1}\sum_{\sigma\in S_n}\sigma$. The pairing
of $x^j$ with $et^j$ is one, and repeated even
constituents contribute their factorials. There is no
additional binomial factor.

For a free jet algebra this identifies the charge piece
with $(K_n\otimes k[x_1,\ldots,x_n])^{S_n}$, where $K_n$
records the finite current and constituent data.
For a jet quotient its dual core is a submodule of this
free polarized module. Multiplication by $d_n$ is
therefore injective. The current differential changes
only how the same constituents are grouped. It leaves
their variables fixed. It is $\Lambda_n$-linear and
extends to \eqref{v4-arith:md:eq:current-localization}.

For two groups of charges $p,q$, the joint recipient
also inverts every cross difference $x_a-y_b$.
For field parameters it inverts their translated
differences as well. At the fixed precision only
finitely many denominator factors occur in any element.
Put $w(u)=1-\hbar/u$ and define
\begin{equation}\label{v4-arith:md:eq:current-fields}
 \begin{gathered}
 \Omega_{p,q}(v)=\prod_{a,b}w(v+x_a-y_b),
 \qquad \mu_w=\mu_0\Omega_{p,q}(0),\\
 Y_B(b,v)c=\mu_w(e^{vT_B}b\otimes c).
 \end{gathered}
\end{equation}
The map $\mu_0$ is ordinary current multiplication
extended to this joint localization. The native vacuum
is the empty symmetric monomial.

A current morphism is an $R_2$-linear unital dg algebra
map of the polynomial cores, linear on the marked
generator spaces. It preserves the scalar-current map
and charge, commutes with $T_B$ and every $D_i$, and
extends continuously to the completed cores.
Its extension to the localizations is then unique.
These conditions preserve the product, fields, exchange
and vacuum in \eqref{v4-arith:md:eq:current-fields}.
The polynomial core is retained as a marking. It need
not be a subalgebra for the deformed product $\mu_w$.

The objects on both sides are scalar extensions from
$k$ to $R_2$. Their morphisms need not be extended from
$k$. After extension, generator duals are $R_2$-linear
with the same topologies. The Frobenius pairing
$(a,b)\mapsto[\hbar](ab)$ identifies these dual calculations
with exact $k$-linear duality.

\begin{proposition}[The two field structures]
\label{v4-arith:md:prop:fields}
The input and current fields just defined satisfy the
differential, vacuum, translation, associativity and
exchange identities on their stated recipients.
Their exchanges are unitary and satisfy Yang--Baxter
and the field hexagon.
\end{proposition}

\begin{proof}
The commuting jet derivations give
$[T\otimes1,C_1(u)]=-\partial_uC_1(u)$ and
$[T\otimes1+1\otimes T,C_1(u)]=0$.
Each derivation kills the vacuum. The product rule
splits a contraction against a product into its two
contractions. Thus either association of three fields
uses the same three pair factors, after translation
of their separations. The reversed ratio of the two
input factors is $\mathcal S_A$.

On the current side, normalized symmetrization absorbs
each intermediate projector. Both associations have
the factor
$\Omega_{p,q}(0)\Omega_{p,r}(0)\Omega_{q,r}(0)$.
The differential leaves constituent variables fixed,
and total translation kills differences. Hence both
are derivations of the deformed product.

Empty groups
give the two vacuum identities. The exchange is
$\Omega_{q,p}(-v)^{21}/\Omega_{p,q}(v)$.
Reversing twice gives its inverse. Commuting pair
factors give Yang--Baxter, and partitioning their pair
sets gives the hexagon. Stability under the jet
derivations carries the input identities through
the quotient ideals.
\end{proof}

\subsection{Formal equivalences and the contact coalgebra}

The completed core has a specified notion of weak
equivalence. Let $S$ be a finite-dimensional nonpositive
augmented commutative dg $k$-algebra with nilpotent
augmentation ideal $\mathfrak m_S$. Put
\[
 \Omega_n=k[t_0,\ldots,t_n,dt_0,\ldots,dt_n]
 /(\textstyle\sum t_i-1,\sum dt_i),
 \qquad d(t_i)=dt_i,
\]
where $|t_i|=0$ and $|dt_i|=1$.
For the completed untwisted core $\widehat{\mathcal C}_B$
with positive symmetric-degree ideal $I_B$, consider
the simplicial set $X_B(S)$.
Its $n$-simplices are continuous unital
$R_2$-linear dg algebra maps satisfying
\begin{equation}\label{v4-arith:md:eq:formal-tests}
 \begin{gathered}
 \phi:\widehat{\mathcal C}_B\longrightarrow R_2\otimes S\otimes\Omega_n,\\
 \phi(I_B)\subset R_2\otimes\mathfrak m_S\otimes\Omega_n.
 \end{gathered}
\end{equation}

In particular $\phi(ab)=\phi(a)\phi(b)$.
The target is discrete. These test maps need not
preserve the generator splitting or the jet markings.
A current morphism $B\to B'$ is a
\emph{formal core equivalence} if precomposition gives
a weak equivalence $X_{B'}(S)\to X_B(S)$ for every $S$.
This defines the weak equivalences on the current side.

The generator images identify
\eqref{v4-arith:md:eq:formal-tests} with the Maurer--Cartan
simplicial set of
$(\mathfrak g\otimes A_{R_2})\otimes\mathfrak m_S$:
the chain-map equation is $d\alpha+\tfrac12[\alpha,\alpha]=0$.
Nilpotence makes all products finite. Continuity gives
finite weight support at each simplex.
The nilpotent Maurer--Cartan invariance theorem sends a
state quasi-isomorphism to an equivalence of these tests.
See Getzler, \emph{Lie theory for nilpotent
$L_\infty$-algebras}, Theorem 4.8.

Conversely, the test $S=k[\varepsilon]/(\varepsilon^2)$
has, at the zero point,
\[
 \pi_nX_B(S)=H^{1-n}(\mathfrak g\otimes A_{R_2})\varepsilon
 \quad(n\geq1),\qquad\pi_0X_B(S)=*.
\]
The bracket is zero in this test, and the polynomial
de Rham complex computes the shifted linear complex.
Every state degree occurs because the states are
nonpositive. Thus a formal core equivalence detects
the recovered state quasi-isomorphism. This uses the
marked core and does not replace it by ordinary
cohomology of the unmarked meromorphic algebra.

Now form the full contact coalgebra $C_A$ of the input
fields. Its singleton is $(\omega_X\otimes A_{R_2})[1]$
with right action \eqref{v4-arith:md:eq:right-current}.
For every ordered partition retain its open and closed
direct images and all normal derivatives. A merger
uses the principal part of the field, together with
exchange to bring nonconsecutive blocks together.

In particular it retains
\begin{equation}\label{v4-arith:md:eq:contact-mark}
 r_A(u^m;a,b)=\operatorname{pp}(u^mY_A(a,u)b)
 \qquad(m\in\mathbb Z).
\end{equation}
Here $\operatorname{pp}$ takes the strictly negative
Laurent part. Fields have expansion
$Y_A(a,u)b=\sum_n(a_{(n)}b)u^{-n-1}$, so
\begin{equation}\label{v4-arith:md:eq:mode-recovery}
 a_{(n)}b=[u^{-1}]r_A(u^n;a,b).
\end{equation}
Negative multipliers recover the regular modes as well.
The coefficient differential and all its coherences
come from the full contact construction, not from
the extraction of a single residue.

Define $U_*(A)=\Omega_\star C_A$. If $tc$ is the
degree-one desuspension of a reduced coalgebra element,
its cobar differential is
\begin{equation}\label{v4-arith:md:eq:contact-cobar}
 d(tc)=-t(d_Cc)+
 \sum_{\overline\Delta c=c'\otimes c''}
       (-1)^{|c'|+1}(tc')(tc'').
\end{equation}
Every ordered-subset coproduct is included. The empty
word is the unit of $U_*(A)$.

\begin{proposition}[The bar contact marking]
\label{v4-arith:md:prop:bar-marking}
The identity at empty support and the reduced map
\begin{equation}\label{v4-arith:md:eq:bar-marking}
 \eta_A(c)=\sum_{r\geq1}(s\tau)^{\star r}
                   \overline\Delta^{(r-1)}c,
 \qquad \tau(c)=tc,\quad |s|=-1,
\end{equation}
give a natural equivalence $C_A\to B_\star U_*(A)$
preserving the entire contact diagram.
\end{proposition}

\begin{proof}
Use bar signs $b_1(sa)=-s(da)$ and
$b_2(sa,sb)=(-1)^{|a|}s(ab)$.
In \eqref{v4-arith:md:eq:contact-cobar}, a split has
sign $(-1)^{|c'|+1}$. It cancels the bar multiplication
joining the corresponding adjacent factors. Internal
differentials have the suspension sign in $b_1$.
Coassociativity identifies the iterated cuts, proving
the chain and coalgebra identities.

Use the finite gap filtration of
Theorem~\ref{v4-arith:md:thm:coefficient-koszul}.
A single original entry gives the identity.
A longer word has a contractible first gap.
On $I$ the filtration ends at $|I|$, proving equivalence.
Every map is a composite of the full coefficient
maps, so it retains collisions, spectators and normal
derivatives. Naturality gives compatibility with
input morphisms.
\end{proof}

Write $\mathcal N_*^{\mathrm{bc}}$ for the input category
with this bar contact marking, and
$\mathcal N_\sharp^{\mathrm{cj}}$ for the marked current
category just defined. The bar marking realizes the
field operations inside $B_\star U_*A$. It is not an
identification with bare multiplication on a singleton
of the cobar algebra.

For $J=C_k$, the state unit and augmentation give
$v_A:J\to C_A$ and $a_A:C_A\to J$, with $a_Av_A=1$.
The bar pointing is $\eta_Av_A:J\to B_\star U_*A$.
It is distinct from the empty-support tensor unit $K$.

\subsection{Native multiplication and every collision}

The current product has a tensor comparison determined
by its own pair multipliers. In the ambient current
category, the tensor of the cores of $A$ and $A'$ is
the core of $A\times A'$. Its localization includes
the cross differences between the two groups.

Let $F_{B,C}=\Omega_{B,C}(0)$ and let $\tau$ be the
Koszul permutation. The deformed crossing and product
on this joint tensor are
\[
 \sigma_w=\tau(F_{C,B}^{21})^{-1}F_{B,C},\qquad
 \mu_{B\otimes_wC}=(\mu_B\otimes\mu_C)
                         (1\otimes\sigma_w\otimes1).
\]
Removing the native cross multiplier has tensor
comparison $U_2=F_{B,C}^{-1}$.
For three groups, both tensor comparisons contain
the same three commuting cross factors.

This proves
their associativity. An empty group gives the unit.
The charge-zero-$k$ slice is not closed under this
tensor, since $(A\times A')_0=k\times k$.
All tensor assertions here use the ambient category.

Let $C_A^0$ denote the contact coalgebra for the
regular fields $\mu_A(e^{uT}\otimes1)$, and write
$C_A^w=C_A$ for the deformed one. On an ordered
configuration multiply all pair operators $E(u)$
from \eqref{v4-arith:md:eq:charged-fields}.
Call this operator $E_A$.

Translation covariance
makes it horizontal, so it extends to closed direct
images and all normal derivatives. The product rule
at a merger gives, before taking principal parts,
\[
 \mu_{ab}e^{uT_a}
 E_{ac}(z_b+u-z_c)E_{bc}(z_b-z_c)
 =E_{(ab)c}(z_b-z_c)\mu_{ab}e^{uT_a}.
\]
The internal factor $E_{ab}(u)$ is precisely the
input contraction. This proves compatibility with
adjacent mergers and their Laurent multipliers.

For a nonconsecutive merger, exchange supplies the
passage through the intervening block. For $13$,
put $u=z_1-z_3$, $t=z_2-z_3$ and
$X_{12}(u-t)=E_{12}(u-t)E_{21}(t-u)^{-1}$.
The corresponding equality is
\begin{align*}
 r^0_{13}E_{12}(u-t)E_{13}(u)E_{23}(t)
 ={}&E_{2,(13)}(t)\operatorname{pp}_{u=0}
       \mu_{13}e^{uT_1}E_{13}(u)X_{12}(u-t).
\end{align*}
It remains valid after every Laurent multiplier.
Commuting the pair factors proves the higher
collision comparisons. Thus $E_A:C_A^w\to C_A^0$
is an isomorphism of full coefficient complexes.

For a cut $G\sqcup H$, let $W_{G,H}$ be the product
of the inverse input contractions between the two
groups. Cancellation of the internal pair factors
gives
\begin{equation}\label{v4-arith:md:eq:cross-cut}
 (E_G\boxtimes E_H)\Delta^w_{G,H}E_A^{-1}
 =W_{G,H}\Delta^0_{G,H}.
\end{equation}
Every ordered-subset cut occurs. For three groups,
both composites contain $W_{G,H}W_{G,J}W_{H,J}$.
These are identities of the full diagrams, including
spectators and normal derivatives.

Recover the state algebra $A_B$ from a current core
and put $G(B)=(C_{A_B}^0)^\sharp$.
This constructs the ordinary contact system before
coefficient duality is applied.
Colored partitions identify the contact system of
$A\times A'$ with the two-color coefficient tensor.

Dualizing this comparison and the native map $U_2$
gives $J^0$. Define its deformed tensor comparison by
$J^w=J^0(W_{G,H})^\sharp$ on each cut component.
The pair factor is forced by expanding the native
multiplier $w(v+x_a-y_b)$ in the constituent jets.
Its first-order expansion is the negative transpose
of $C_1(v)$, so it is the same $W_{G,H}$ as in
\eqref{v4-arith:md:eq:cross-cut}.

Consequently actual native multiplication has the
comparison
\begin{equation}\label{v4-arith:md:eq:native-product-square}
 G(\mu_w)J^w
 = (\Delta^0_{G,H})^\sharp(W_{G,H})^\sharp.
\end{equation}
Indeed removing the cross multiplier turns native
multiplication into ordinary core multiplication.
The latter transposes the diagonal $A\to A\times A$,
whose colored-partition map is the cut coproduct.
Restoring the native multiplier gives the right
side of \eqref{v4-arith:md:eq:native-product-square}.
Equation~\eqref{v4-arith:md:eq:cross-cut} now shows that
$E_A^\sharp$ identifies this realized algebra with
$(C_A^w)^\sharp$. Denote the resulting operation-preserving
realization by $U_\sharp$.

\begin{theorem}[Bar-marked charged-current duality]
\label{v4-arith:md:thm:native}
At precision $R_2$ there is an anti-equivalence
\[
 D:(\mathcal N_*^{\mathrm{bc}})^{\mathrm{op}}
 \xrightarrow{\sim}\mathcal N_\sharp^{\mathrm{cj}},
 \qquad D(A)=B_A^{\mathrm{mer}}.
\]
Its inverse recovers the $R_2$-algebra
\[
 R(B)=\operatorname{Hom}_{\mathfrak g_{R_2}}
       (\mathfrak g_{R_2},W_B^\vee[-1]).
\]
There is a natural coefficient equivalence
\begin{equation}\label{v4-arith:md:eq:native-square}
 \mathcal K(U_*A)
 =(B_\star\Omega_\star C_A)^\sharp
 \xrightarrow{\eta_A^\sharp}C_A^\sharp
 \xrightarrow{\sim}U_\sharp D(A).
\end{equation}
It preserves the full operations and both native
vacuum identities. It takes the bar pointing to
its outgoing transpose to $J^\sharp$.
The comparison remains an anti-equivalence after
the declared weak equivalences are inverted,
including their mapping spaces.
\end{theorem}

\begin{proof}
The adjoint representation of $\mathfrak{sl}_2$ has
\[
 \operatorname{End}_{\mathfrak g}(\mathfrak g)=k,
 \qquad\operatorname{Hom}_{\mathfrak g}
       (\mathfrak g\otimes\mathfrak g,\mathfrak g)=k[\,,\,].
\]
Its exterior square is the adjoint representation.
Its symmetric square is the sum of the five-dimensional
irreducible and the invariant line.

Thus
\eqref{v4-arith:md:eq:isotypic} forces the recovered
bracket to be $[x,y]\otimes m(a,b)$.
Graded antisymmetry makes $m$ graded commutative.
Jacobi on $e\otimes a,e\otimes b,f\otimes c$ gives
$a(bc)=(-1)^{|a||b|}b(ac)$, and hence
\[
 (ab)c=(-1)^{(|a|+|b|)|c|}c(ab)
       =(-1)^{|b||c|}a(cb)=a(bc).
\]
The scalar current supplies the unit.

Equivariance makes the differential, translation and
each marked jet derivation equal to $1\otimes d_A$,
$1\otimes T$ and $1\otimes\delta_i$.
The Lie derivation identities become the product
derivation identities. A transposed core morphism
is $1\otimes f$. Its bracket equation is exactly
multiplicativity of $f$. Its other equations give
the stated input morphism conditions.
This recovers objects and maps.

The inverse comparisons are $a\mapsto(x\mapsto
x\otimes a)$ and inverse dual evaluation of
$\mathfrak g\otimes R(B)\to L_B$.
They extend uniquely through the completed core
and its localization. They commute with the
symmetric-ring actions and pair multipliers, so
preserve the actual native fields and products.
Their triangles hold on generators and therefore
on the complete presentations.

The formal-test calculation identifies the two
classes of weak equivalences. The exact $R_2$-dual
is justified by its Frobenius pairing. Thus the
inverse functors descend to the localized categories.
The universal property of localization retains the
inverse comparisons and all mapping spaces.

Finally \eqref{v4-arith:md:eq:native-product-square}
and the bar marking give
\eqref{v4-arith:md:eq:native-square}. The three-group
factor identity proves coherence with the tensor
comparisons. The ambient field unit, obtained from
the zero state algebra, is $R_2$ and realizes to
$K$. Its unit map supplies the native vacuum.

The scalar-current object $B_k$ instead realizes
to $J^\sharp$. The map to that object is
$v_A^\sharp$, and
$v_A^\sharp\eta_A^\sharp=(\eta_Av_A)^\sharp$.
The augmentation supplies its section. This proves
the pointing and both vacuum assertions.
\end{proof}

The theorem uses the bar contact marking, all jet
actions and the split quadratic current core.
An arbitrary enhancement with the same unmarked
coefficient algebra is a larger recognition problem.
Even inside the stated categories, coefficient
framings and coefficient arrows contain information
that an unframed state isomorphism forgets.

\section{The cubic enhancement space and its arrows}
\label{v4-arith:md:sec:enhancements}

An enhancement of a coefficient algebra $P$ is a
native object $A$ together with an equivalence
$U(A)\simeq P$. Work in the categories localized
at their declared weak equivalences. The space of
enhancements is the homotopy fiber, over $P$, of
the map between the maximal subgroupoids of
the native and coefficient categories.

Take $A_0$ from \eqref{v4-arith:md:eq:cubic}, and put
$C_0=C_{A_0}$, $P_0=\Omega_\star C_0$ and
$E_0=C_0^\sharp$. We write coefficient maps in
coalgebra coordinates on $C_0$. Bar--cobar identifies
them with maps of $P_0$. Transposition identifies
them, with composition reversed, with maps of $E_0$.
A scalar-pointed map commutes with both maps induced
by $k\to A_0\to k$.

\subsection{All coefficient endomorphisms}

Let $\mathfrak m=V\oplus Q$. For $v\in A_0[z]$ and
$p\in k[z]$, define
\[
 J_n(v)=\frac{(\partial_z+T)^n}{n!}v,
 \qquad
 S_p(e_n)=J_n(pe_0)
 =\sum_{j=0}^n\frac{p^{(j)}(z)}{j!}e_{n-j}.
\]
Extend $S_p$ multiplicatively, fixing $k[z]$.
Horizontality of a singleton map $F(z)$ means
$[\partial_z+T,F]=0$. Its field equation is
\begin{equation}\label{v4-arith:md:eq:geometric-field-map}
 F(y)Y(a,u)b=Y(F(y+u)a,u)F(y)b.
\end{equation}
The shifted position of the first input is essential.

For this full contact system, a coefficient map is
determined by its singleton. The residue-lattice
description proves this. On $n$ positions the
contact complex is a module in degree $-n$.
It embeds in the open direct image from the fully
distinct locus.

The module is the intersection of the pair
residue kernels. Cuts determine a map on that locus.
The embedding then determines its extension.
Maps between equally shifted modules have no
negative Ext groups, so their mapping spaces are
discrete. Limits over the collision diagram retain
this discreteness.

For completeness, concentration of the contact
complex follows by induction on the number of
positions. Off the small diagonal it splits into
two smaller groups. The remaining cokernel is a
principal-part residue on the small diagonal.
Insert the vacuum in one colliding input to lift
each principal part.

Its other residues vanish by
the residue-kernel condition. This proves
surjectivity and concentration. The finite complex
of flat coefficients remains exact under coefficient
extension. This also justifies its use over $R_2$.

Naturality of every Laurent-multiplied binary
boundary gives \eqref{v4-arith:md:eq:geometric-field-map}.
Both cut orders retain unmultiplied exchange.
Conversely a horizontal singleton satisfying these
field and exchange equations is placed at each
block. The field equation intertwines adjacent
residues, exchange intertwines nonconsecutive
passages, and horizontality retains normal
derivatives. It therefore extends to the full
coefficient map.

\begin{theorem}[Coefficient endomorphisms of the cubic system]
\label{v4-arith:md:thm:coefficient-endomorphisms}
Every scalar-pointed coefficient endomorphism is
uniquely $F_{p,v}$, with $p\in k[z]$ and
$v\in\mathfrak m[z]$, where
\begin{equation}\label{v4-arith:md:eq:coefficient-endomorphisms}
 F_{p,v}(1)=1,
 \qquad F_{p,v}(e_n)=J_n(pe_0)+\hbar J_n(v),
\end{equation}
extended multiplicatively. Composition is
\begin{equation}\label{v4-arith:md:eq:coefficient-composition}
 (p,v)\circ(q,w)=(pq,qv+S_p(w)).
\end{equation}
\end{theorem}

\begin{proof}
Reduce $F$ modulo $\hbar$. The constant field
coefficient makes $F_0$ an augmented algebra map.
Let $f=F_0(e_0)$. If its largest jet index is
$N>0$, then
$\delta_{N+1}f=e_0\partial f/\partial e_N\ne0$.
In charge at most two, multiplication by $e_0$
does not kill this nonzero derivative.

The exchange equation on $e_0\otimes e_0$ is
\[
 C_1(u)(f(y+u)\otimes f(y))
 =-u^{-1}f(y+u)\otimes f(y).
\]
The left side has nonzero coefficient
\[
 [u^{-2N-1}]=(-1)^{N+1}\binom{2N}{N}
       \delta_{N+1}f(y)\otimes\delta_{N+1}f(y).
\]
The two tensor factors are free over $k[y]$.
The right side has at most a simple pole, a
contradiction. Hence $N=0$.

Write $f=pe_0+be_0^2$. The simple-pole equation
is $\delta_1f\otimes\delta_1f=f\otimes f$.
Its bidegree $(2,2)$ coefficient gives
$4b^2=b^2$, hence $b=0$. Horizontality then
forces $F_0=S_p$. The displayed formula for
$S_p$ shows that it commutes with every jet
derivation and satisfies the field equation.

Write $F=S_p+\hbar G$. A first-order correction
and a quantum contraction together have order
$\hbar^2$, so the remaining equation is classical.
It says $G(1)=0$ and
$G(ab)=G(a)S_p(b)+S_p(a)G(b)$.

Horizontality determines $G(e_n)=J_n(v)$ from
$v=G(e_0)\in\mathfrak m[z]$. Such a derivation
preserves $I^3$, since each generator image has
positive charge. The singleton criterion proves
existence and uniqueness of the full map.
Evaluating a composite on $e_0$ gives
\eqref{v4-arith:md:eq:coefficient-composition}.
\end{proof}

For example $S_z(e_n)=ze_n+e_{n-1}$, where
$e_{-1}=0$. The term $e_0$ in $S_z(e_1)$ is
forced by horizontality. This map commutes with
all jets but is not a constant native state map.

If only the empty unit and counit are fixed,
there is also a map $Z$ vanishing on every
nonempty coefficient. Every other map is locally
unital and has \eqref{v4-arith:md:eq:coefficient-endomorphisms},
with $v\in A_0[z]$ permitted when $p=0$.
Indeed the field equation makes $F(1)$ an
idempotent, hence zero or one. A zero value
forces the map to vanish on nonempty support.

For value one and $p\ne0$, a scalar part
$v_{\mathrm{sc}}$ would give
$3\hbar p^2v_{\mathrm{sc}}e_0^2$ in the cubic
relation, so it must vanish.
For $p=0$, products of two positive generator
images already vanish, and impose no such condition.
Composition still has the displayed formula,
and $Z$ is absorbing. In particular
$F_{0,v}F_{0,w}=F_{0,w_{\mathrm{sc}}}$.
The scalar-pointed local projection $F_{0,0}$
preserves the local vacuum and differs from $Z$.

\subsection{The entire enhancement fiber}

An invertible coefficient map has $p=a\in k^\times$.
In the coordinates $F_{a,av}$ its group is
\begin{equation}\label{v4-arith:md:eq:framing-group}
 \Gamma=\mathfrak m[z]\rtimes k^\times,
 \qquad(v,a)(w,b)=(v+\vartheta_aw,ab),
 \quad\vartheta_a(w_V+w_Q)=w_V+aw_Q.
\end{equation}
The native endomorphisms are exactly
$s_\lambda(e_n)=\lambda e_n$, $\lambda\in R_2$.
Charge and $\ker(\delta_2|_V)=ke_0$ determine
the seed image, and translation and multiplication
determine the map. The native automorphism image
is therefore
$H=\{(be_0,a):b\in k,\ a\in k^\times\}$.

\begin{theorem}[Cubic enhancement fibers]
\label{v4-arith:md:thm:fibres}
Within the two stated native categories, every
enhancement of $P_0$, respectively $E_0$, is
equivalent to $A_0$, respectively $D(A_0)$,
with a coefficient framing. Its enhancement
space is
\begin{equation}\label{v4-arith:md:eq:fibres}
 \operatorname{Enh}_*(P_0)
 \simeq\operatorname{Enh}_\sharp(E_0)
 \simeq(\Gamma/H)_{\mathrm{disc}}
 \simeq(\mathfrak m[z]/ke_0)_{\mathrm{disc}}.
\end{equation}
The quotient $\Gamma/H$ is a right-coset set.
\end{theorem}

\begin{proof}
Bar--cobar turns an enhancement into $C_A\simeq C_0$.
Singleton cohomology and faithful scalar extension
by $k[z]$ and $R_2$ force $H^i(A)=0$ for $i\ne0$.
The map $A\to H^0(A)$ is a charged jet
quasi-isomorphism: its boundary ideal is stable
under all markings. We may therefore work with
a degree-zero algebra.

The residue-lattice criterion gives an invertible
horizontal polynomial singleton. Modulo $\hbar$,
its constant field coefficient is an algebra
isomorphism over $k[z]$. Evaluation at $z=0$
gives $A\cong A_0$ as ordinary algebras.

If $d=\dim E_A$, its map on augmentation
indecomposables identifies $\mathcal D_X^d$
with $\mathcal D_X$. These are the free
Weyl modules generated by the seeds.
For the Bernstein filtration,
$\dim_kF_m\mathcal D_X=(m+1)(m+2)/2$.
An isomorphism and its inverse have bounded
filtration displacement, so comparison of the
leading quadratic growth gives $d=1$.

The ordinary algebra isomorphism transports
$\mathfrak m^3=0$ and the isomorphism
$\operatorname{Sym}^2(\mathfrak m/\mathfrak m^2)
\to\mathfrak m^2$. Generation in charge one
therefore leaves no quadratic relation and
kills every cubic product. The one-dimensional
seed gives precisely the ideal $I^3$.
Thus $A$ is a charged jet copy of $A_0$.

Maps between degree-zero objects in this localized
nonpositive dg category are their ordinary discrete
sets. Indeed $H^0$ is left adjoint to the inclusion
of degree-zero objects. It preserves weak
equivalences and has identity counit there.
Coefficient maps are discrete by the lattice
criterion.

The enhancement fiber is consequently
the right action groupoid $\Gamma//H$.
This action is free, so the groupoid is equivalent
to the discrete coset set. Right multiplication
by $(be_0,c)$ sends $(v,a)$ to $(v+be_0,ac)$,
giving the last identification in
\eqref{v4-arith:md:eq:fibres}.
The native duality proves the same assertion
on the opposite side.
\end{proof}

The set in \eqref{v4-arith:md:eq:fibres} is not
a quotient group. For $a\ne1$,
\[
 (e_0^2,1)(0,a)(e_0^2,1)^{-1}
 =((1-a)e_0^2,a)\notin H.
\]
The transitive action groupoid of $\Gamma$ on
its coset set is $BH$. It differs from the
enhancement fiber, whose action by $H$ on
$\Gamma$ is free.

\subsection{Relative arrows and a higher-precision defect}

For a realization $U:\mathcal N\to\mathcal A$ and a
coefficient chain $a_0\to\cdots\to a_n$, take the
homotopy fiber of
\[
 \operatorname{Fun}([n],\mathcal N)^\simeq
 \longrightarrow\operatorname{Fun}([n],\mathcal A)^\simeq
\]
over that chain. The superscript means the maximal
subgroupoid. This \emph{relative chain space} retains
the enhancements at all vertices and their compatible
native arrows. The case $n=0$ is the enhancement
space above.

\begin{proposition}[Lifting coefficient arrows]
\label{v4-arith:md:prop:arrows}
For fixed frames $g_s,g_t\in\Gamma$, a coefficient
map $F$ has a native lift exactly when
\begin{equation}\label{v4-arith:md:eq:lift}
 F=g_t s_\lambda g_s^{-1}
 \qquad\text{for some }\lambda\in R_2.
\end{equation}
The lifting space is empty or a point.
Allowing both frames to vary, $F_{p,v}$ lifts
somewhere precisely when $p\in k^\times$, or
$p=0$ and $v\in ke_0$. The map $Z$ never lifts.
\end{proposition}

\begin{proof}
Theorem~\ref{v4-arith:md:thm:fibres} supplies all
objects, and the maps $s_\lambda$ are all their
native endomorphisms. Distinct $\lambda$ have
distinct singleton values, which proves the
fixed-frame assertion. If $\lambda$ has nonzero
reduction it is a unit, and changing one frame
produces every coefficient automorphism.

For $\lambda=\hbar b$, let $a_s,a_t$ be the
classical scalings of the two frames. The
composition law gives
\[
 g_t s_{\hbar b}g_s^{-1}
 =s_{\hbar b a_t/a_s}.
\]
Every first-order correction in a frame acquires
a second factor of $\hbar$ and vanishes.
Only the constant $e_0$ direction can therefore
lift in this case. Nonconstant $p$ and the map
$Z$ meet neither possibility.
\end{proof}

For a chain $(F_i)$, all lifts are consequently
described by $F_i=g_i s_{\lambda_i}g_{i-1}^{-1}$,
modulo changes at each vertex
\[
 g_i\mapsto g_i h_i^{-1},\qquad
 s_{\lambda_i}\mapsto h_i s_{\lambda_i}h_{i-1}^{-1},
 \qquad h_i\in H.
\]
The action on the frames is free. This gives a
discrete relative chain space.

Faces compose
the adjacent scalars and degeneracies insert $1$.
The anti-equivalences and
\eqref{v4-arith:md:eq:native-square} identify it
with the relative chain space over
$\mathcal K(a_n)\to\cdots\to\mathcal K(a_0)$.
Indeed a frame $g$ becomes $(g^{-1})^\sharp$, and
\[
 (g_s^{-1})^\sharp s_\lambda^\sharp g_t^\sharp
 =(g_t s_\lambda g_s^{-1})^\sharp.
\]
This identity retains empty lifting spaces as well
as the nonempty ones.

For instance $S_z$ has no lift under any reframing.
Neither does $F_{0,e_1}$, whose singleton satisfies
\[
 F_{0,e_1}(e_n)=\hbar(n+1)e_{n+1},
 \qquad[\delta_2,F_{0,e_1}](e_0)=\hbar e_0.
\]
The latter preserves the coefficient field diagram
but not the marked native jet action. Its relative
space is empty on both sides of the duality.

The fixed first-order precision is also substantive.
Define a horizontal derivation on $A_0$ by
\[
 D(e_n)=\sum_{i+j=n}e_ie_j.
\]
The formula satisfies $[D,T]=0$, as follows by
differentiating the two factors and collecting
their coefficients. The map $1+\hbar D$ is
a coefficient automorphism over $R_2$.

\begin{proposition}[A frame obstructed at the next precision]
\label{v4-arith:md:prop:precision}
Keep the exchange kernel $w(u)=1-\hbar/u$.
The frame $1+\hbar D$ has no extension over
$R_3=k[\hbar]/(\hbar^3)$ preserving exchange,
even with an arbitrary position-dependent
second-order correction.
\end{proposition}

\begin{proof}
Write a proposed extension as
$1+\hbar D+\hbar^2K$. Exchange has expansion
$1+2\hbar C_1+\hbar^2\mathcal S_2$.
Its order-two intertwining defect is forced to be
$2[C_1,D\otimes1+1\otimes D]$.
Both the new correction $K$ and $\mathcal S_2$
commute with the identity at this order.

On $e_0\otimes e_0$ the defect is
\begin{equation}\label{v4-arith:md:eq:precision-defect}
 -\frac2u(e_0^2\otimes e_0+e_0\otimes e_0^2)\ne0.
\end{equation}
The terms have separate bidegrees $(2,1)$ and
$(1,2)$. Multiplying the two tensor factors
would kill both by $I^3=0$, but exchange acts
before that multiplication. No second-order
correction can remove the displayed defect.
\end{proof}

This is an obstruction to extending the specified
frame with the fixed exchange kernel. The native
duality and the exhaustive fiber calculation
remain assertions at $R_2$. A compatible tower
requires its own objects, surviving frames and
comparison maps at every precision. More general
native enhancements also require a separate
recognition theorem or an obstruction on their
declared domain.

\section{Fields can forget a jet marking}
\label{v4-arith:md:sec:filtered-moments}

The constituent identities in the charged-current
presentation contain more information than the
commutation relations of its jet derivations.
We now vary those identities while retaining the
entire rational field. This gives a larger
marked category in which the standard recognition
statement fails, together with a duality on an
explicit family inside that category.

Continue over $R_2=\mathbb C[\hbar]/(\hbar^2)$.
Use the cubic algebra
$A=R_2\oplus V_{R_2}\oplus\operatorname{Sym}^2V_{R_2}$
of \eqref{v4-arith:md:eq:cubic}, with
$V=\bigoplus_{n\geq0}\mathbb Ce_n$.
The differential is zero, $e_n$ has charge one
and weight $n+1$, and $\hbar$ has weight zero.

Keep $T$ and the standard $\delta_i$ of
\eqref{v4-arith:md:eq:jets}, and set $N=\delta_1$.
On charge one,
\[
 \delta_{m+1}=\delta_2^m.
\]
This identity concerns compositions on $V$.
It does not identify $\delta_{m+1}$ with a
power of a derivation on the whole algebra.

\begin{definition}[Filtered jet marking]
\label{v4-arith:md:def:filtered-marking}
A \emph{filtered jet marking} on a nonpositive
commutative dg state algebra consists of degree-zero
chain derivations $T,J_i$ satisfying
\[
 J_1=N,\qquad [J_i,J_j]=0,\qquad
 [J_i,T]=(i-1)J_{i-1},\qquad J_0=0.
\]
They annihilate the unit and preserve charge.
Each $J_i$ is bounded for the weight filtration.

Only finitely many act nontrivially on each
state. The rational fields and exchange are
\eqref{v4-arith:md:eq:charged-fields} with
$\delta_i$ replaced by $J_i$.
Morphisms are constant $R_2$-linear dg algebra
maps preserving charge, translation, every
$J_i$, the fields and exchange.
Weak equivalences are the state
quasi-isomorphisms among these maps.
\end{definition}

In particular $[d,J_i]=0$ is required.
Exact operator weight $1-i$ and the constituent
identities on charge one are additional conditions.
The following family satisfies neither additional
condition when its parameter is nonzero.

\begin{proposition}[Equal fields with distinct markings]
\label{v4-arith:md:prop:moment-family}
For $\kappa\in\mathbb C$, put
\[
 J_1^\kappa=N,\qquad J_2^\kappa=\delta_2,\qquad
 J_i^\kappa=\delta_i+
       \hbar\kappa\binom{i-1}{2}\delta_{i-2}
       \quad(i\geq3).
\]
These are filtered jet markings on $A$.
Write $A^\kappa$ for the resulting marked algebra.
Their rational fields, exchange and full
collision coefficients are identical.
Distinct parameters remain inequivalent after
inverting marked state quasi-isomorphisms.
For $\kappa\ne0$ they are not equivalent there
to a standard charged-jet quotient.
\end{proposition}

\begin{proof}
The operators are linear combinations of commuting
derivations. Translation covariance follows from
\[
 \binom{i-1}{2}(i-3)=(i-1)\binom{i-2}{2}
\]
for $i\geq4$. The added term at $i=3$ is
$\hbar\kappa N$, which commutes with $T$.
On $e_n$ both terms vanish for $i>n+3$.
The derivation rule proves finite support on
every state. Their two weight shifts are
$1-i$ and $3-i$.

The contraction in
\eqref{v4-arith:md:eq:charged-fields} has an
external factor $\hbar$. Each change of a
marked operator has another factor $\hbar$.
Thus $\hbar C_1^{J^\kappa}=\hbar C_1^\delta$
over the common base $R_2$. The fields agree
before principal parts are taken.
Every iterate, normal derivative, spectator
and collision comparison therefore agrees.

On the charge-one complex define
\[
 \Xi_2=J_3-J_2^2.
\]
It acts on cohomology because the markings
are chain maps. In this family
\begin{equation}\label{v4-arith:md:eq:moment-defect}
 \Xi_2^\kappa=\hbar\kappa\,1_{V_{R_2}}.
\end{equation}

An invertible intertwiner between parameters
$\kappa,\lambda$ would imply
$\hbar(\kappa-\lambda)1_{V_{R_2}}=0$.
Freeness forces $\kappa=\lambda$.
The same reasoning applies to every zigzag,
since its cohomology maps are invertible
intertwiners. A standard charged-jet quotient
has $\Xi_2=0$, which gives the last assertion.
\end{proof}

The smallest calculation is $J_2e_0=0$ and
$J_3^1e_0=\hbar e_0$. The interaction is also
nonzero:
\[
 [\hbar u^{-1}]Y(e_0,u)e_0=e_0^2.
\]
The quotient by the square of the augmentation
ideal would retain a marking defect but kill
this product. Nor can the infinite jet orbit
be replaced by a nonzero finite-dimensional
charge-one space in characteristic zero:
$[J_2,T]=1$ would contradict the trace of a
commutator.

\subsection{The chain condition detects the distinction}

If the marking need not commute with the
differential, a contractible extension can
erase the preceding distinction.
Let $\widetilde A$ be the graded commutative
algebra on species $e_n,v_n,u_n$, truncated
at total charge three. Give $e_n,v_n$ degree
zero and $u_n$ degree $-1$, with
$du_n=v_n$ and $de_n=dv_n=0$.
All species have weight $n+1$.
Use standard shifts and translation on each.

On their linear span set
$Me_n=0$, $Mv_n=e_n+v_n$, $Mu_n=u_n$.
Set $\widetilde J_1=N$ and
$\widetilde J_2=\delta_2$. For $i\geq3$
define on generators
\[
 \widetilde J_i x_n=x_{n-i+1}
   +\hbar\binom{i-1}{2}Mx_{n-i+3},
\]
then extend as derivations.
The shifts commute with $M$, so these
derivations commute and obey the same
translation identity. Their fields are still
the standard dg fields: a correction to a
marking acquires a second $\hbar$ in the
contraction.

There are marked dg algebra maps
\[
 A^0\xrightarrow{i_0}\widetilde A
 \xleftarrow{i_1}A^1,\qquad
 i_0(e_n)=e_n,\quad i_1(e_n)=e_n+v_n.
\]
To verify their quasi-isomorphism property,
let $p$ kill $u_n,v_n$ and fix $e_n$.
Let $K$ be the degree-$-1$ derivation
$K(v_n)=u_n$, $K(e_n)=K(u_n)=0$.
Then $dK+Kd$ counts the factors from the
two contractible species.

On a monomial containing $r>0$ such factors
put $H=K/r$, and put $H=0$ on pure
$e$-monomials. Here $r$ is one or two.
Consequently
\[
 dH+Hd=1-i_0p.
\]
This proves that $i_0,p$ are quasi-isomorphisms.
Since $pi_1=1$, so is $i_1$.

The marking nevertheless fails the chain
condition:
\begin{equation}\label{v4-arith:md:eq:chain-defect}
 [d,\widetilde J_3]u_0
 =\hbar v_0-\hbar(e_0+v_0)=-\hbar e_0.
\end{equation}
The cohomologous representatives $e_0$
and $e_0+v_0$ have images differing by a
nonzero cohomology class. Thus $\widetilde J_3$
does not induce a cohomology operator.
This span is excluded precisely by the
chain condition in
Definition~\ref{v4-arith:md:def:filtered-marking}.

\subsection{All markings with fixed first jets}

\begin{theorem}[Polynomial classification on the cubic carrier]
\label{v4-arith:md:thm:moment-classification}
Fix $A,T,N$ and $J_2=\delta_2$.
Every filtered jet marking reducing to
$\delta_i$ modulo $\hbar$ has the unique form
\begin{equation}\label{v4-arith:md:eq:polynomial-moments}
 \begin{gathered}
 J_{m+1}^c=\delta_{m+1}
  +\hbar\sum_{r=2}^m\binom mr c_r\delta_{m-r+1},\\
 c(t)=\sum_{r\geq2}c_r\frac{t^r}{r!}
          \in t^2\mathbb C[t].
 \end{gathered}
\end{equation}
All these markings have the same fields.
Distinct polynomials give distinct marked
equivalence classes, including under
chain-compatible nonpositive dg zigzags.
Exact operator weight $-m$ for $J_{m+1}$
forces $c=0$.
\end{theorem}

\begin{proof}
Write $D=\delta_2$. Commutation with $D$
puts $J_{m+1}e_0$ in
$\ker(D|_V)=R_2e_0$, say $a_me_0$.
The first two markings give $a_0=1,a_1=0$.
Reduction gives $a_m=\hbar c_m$ for $m\geq2$.
Finite support on $e_0$ makes $(c_m)$
finitely supported.

Translation covariance is the recursion
\[
 (n+1)J_{m+1}e_{n+1}
   =TJ_{m+1}e_n+mJ_me_n.
\]
Starting at $e_0$, induction gives
\[
 J_{m+1}e_n
  =\sum_{q=0}^{\min(m,n)}
       \binom mq a_{m-q}e_{n-q}.
\]
The induction uses
$(n-q+1)\binom mq+m\binom{m-1}{q-1}
=(n+1)\binom mq$.
Substitute the $a_m$ to obtain
\eqref{v4-arith:md:eq:polynomial-moments}.
The derivation property determines the
operator on all products.

Conversely the displayed operators commute.
The identity
$\binom mr(m-r)=m\binom{m-1}r$
proves translation covariance.
If $c_r=0$ for $r>R$, then
$J_{m+1}e_n=0$ for $m>n+R$.
The same external-$\hbar$ calculation as
before proves equality of fields.

A marked equivalence induces an invertible
cohomology map commuting with $D,T$.
It acts on the $D$-kernel by a unit
$a\in R_2^\times$. Intertwining
$J_{m+1}e_0$ gives
$\hbar(c_m-c'_m)a=0$, hence $c_m=c'_m$.
Finally an operator of weight $-m$
sends $e_0$ to weight $1-m$.
There is no such charge-one state for
$m\geq2$, forcing every $c_m$ to vanish.
\end{proof}

This classifies markings on the fixed
cubic carrier, with the specified first
jets and reduction. It does not classify
all algebras realizing its coefficient
diagram. Write $A^c$ for the algebra
with polynomial marking $c$.

\subsection{Compensated moments and the current partner}

The current partner retains the marked
operators without treating them as
constituent power sums. Use the polynomial
Chevalley core of
$\mathfrak{sl}_2\otimes A$ from
Section~\ref{v4-arith:md:sec:charged}.
Its negative transposes are
\[
 T_B=-(1\otimes T)^\vee,\qquad
 D_i^c=-(1\otimes J_i^c)^\vee.
\]

On total charge $n$, let
$p_m=\sum_{a=1}^n x_a^m$, with $p_0=n$.
Constituent polarization gives
\begin{equation}\label{v4-arith:md:eq:compensated-moments}
 \begin{split}
 P_m^c=-D_{m+1}^c
   &=p_m+\hbar\sum_{r=2}^m\binom mr c_rp_{m-r},\\
 \widehat P_m
   &=P_m^c-\hbar\sum_{r=2}^m\binom mr c_rP_{m-r}^c
     =p_m.
 \end{split}
\end{equation}
The second equality uses $\hbar^2=0$.
It supplies the ordinary action of
$\Lambda_n=\mathbb C[x_1,\ldots,x_n]^{S_n}$.

Define $B^c$ by localizing the charge-$n$ core at
$d_n=\prod_{a<b}(x_a-x_b)^2$ using this
compensated action, with $d_0=d_1=1$.
Use the same cross differences and their
declared translates for multiple inputs.
Define the product and fields by the pair
multiplier
\[
 \Omega_{p,q}(v)
  =\prod_{a=1}^p\prod_{b=1}^q
       \left(1-\frac{\hbar}{v+x_a-y_b}\right),
 \qquad \mu_w=\mu_0\Omega_{p,q}(0).
\]

They are independent of $c$.
Each $D_i^c$ is a finite linear combination
of standard current derivations, so it
extends through these localizations and
remains a derivation of $\mu_w$.
The three-input multiplier is again
$\Omega_{p,q}\Omega_{p,r}\Omega_{q,r}$.
Thus the differential, associativity,
exchange and vacuum identities hold on
the same recipients.

A \emph{compensated current presentation}
retains the split core, scalar current,
charge, translation, $D_i^c$, compensated
spectral action and rational recipients.
Its maps preserve these structures and
are linear on core generators.
It has the formal-core weak equivalences
defined in Section~\ref{v4-arith:md:sec:charged}.
Standardization replaces $D_{m+1}^c$ by
$\widehat D_{m+1}=-\widehat P_m$.
It leaves the underlying map, core,
differential, fields and recipients fixed.
Every specified morphism preserves
$\widehat P_m$, so standardization is a
functor to the standard current category.

Apply $U_\sharp$ to this standardized
presentation, using the algebra structure
induced by its native product and tensor
comparison. This defines $U_\sharp$ on
the displayed compensated category.
Standardization is not claimed to be
an equivalence of marked fields or to
make this category closed under tensor
products.

Let $C_0$ be their common full coefficient
coalgebra and $P_0=\Omega_\star C_0$.
The bar marking and the native product
comparison give a natural equivalence
\begin{equation}\label{v4-arith:md:eq:compensated-comparison}
 (B_\star P_0)^\sharp
  \xrightarrow{\eta^\sharp}C_0^\sharp
  \xrightarrow{\theta_c}U_\sharp(B^c).
\end{equation}
Here $\theta_c$ is the inverse of the
standard current comparison after
standardization. It preserves all
collisions, cuts, normal derivatives,
spectators and local pointings.
The bar pointing remains the composite
$J\to C_0\to B_\star P_0$, separate from
the empty-support unit.

Raw moments would give the wrong
localization. In the one-parameter
family and charge one,
\[
 P_2^\kappa-(P_1^\kappa)^2
       =\hbar\kappa\,1,\qquad
 D_3^\kappa+(D_2^\kappa)^2
       =-\hbar\kappa\,1.
\]
Since $p_2=p_1^2$ in $\Lambda_1=\mathbb C[x]$,
the raw assignment $p_m\mapsto P_m^\kappa$
does not define an action when $\kappa\ne0$.
Compensation repairs this equation while
retaining the marked operator invariant.

\begin{proposition}[All arrows in the displayed family]
\label{v4-arith:md:prop:moment-arrows}
The native maps $A^c\to A^d$ are exactly
\[
 \begin{gathered}
 s_\alpha(1)=1,\quad s_\alpha(e_n)=\alpha e_n,\quad
 s_\alpha(e_ie_j)=\alpha^2e_ie_j,\\
 \hbar\alpha(c_r-d_r)=0\quad(r\geq2),
 \end{gathered}
\]
with $\alpha\in R_2$.
Current pullback gives an anti-equivalence
between the displayed state and compensated
current categories. It commutes with
\eqref{v4-arith:md:eq:compensated-comparison}.
\end{proposition}

\begin{proof}
Charge and $D$ preservation put the image
of $e_0$ in $R_2e_0$.
Translation determines every $e_n$,
and multiplication determines the
quadratic states.
Intertwining $J_{r+1}$ on $e_0$ gives
the condition on $\alpha$.
Formula~\eqref{v4-arith:md:eq:polynomial-moments}
makes it sufficient on all generators,
and product defects acquire further
powers of $\alpha$.

For the reverse construction, scalar-current
equivariance makes a transposed linear
core map $1\otimes f$.
The quadratic Chevalley equation says
$f(ab)=f(a)f(b)$. The scalar-current
condition gives the unit.
Translation and the marked derivations
recover the remaining input equations.
This recognition calculation does not
use the higher constituent identities.

Pullback by $s_\alpha$ scales charge $n$
by $\alpha^n$. It commutes with the
compensated spectral action and extends
through the same localizations.
If $c\ne d$, the condition forces
$\alpha\in\hbar\mathbb C$.

Charge-one images then have a factor
$\hbar$, and all images of charge at
least two vanish. The residual marked
differences therefore vanish on the
image, without inverting a zero divisor.
The native unit and counit are inverse
evaluation on the adjoint-isotypic core,
as in Theorem~\ref{v4-arith:md:thm:native}.
Their triangles and the standardized
comparison prove the assertions.
\end{proof}

For equal parameters every $\alpha$ occurs.
For unequal parameters exactly
$\alpha\in\hbar\mathbb C$ occurs.
Two arrows between unequal parameters
compose to $s_0$, the scalar projection
followed by its section. This map
preserves the native unit.
Isomorphisms require equal parameters
and $\alpha\in R_2^\times$.

The coefficient frame group $\Gamma$
and subgroup $H$ of
\eqref{v4-arith:md:eq:framing-group}
are unchanged. Thus the framed subspace
supported on these displayed classes is
\[
 \coprod_{c(t)\in t^2\mathbb C[t]}
             (\Gamma/H)_{\mathrm{disc}}.
\]
For frames $g_s,g_t$, a coefficient arrow
$F$ has a native lift between parameters
$c,d$ exactly when
\[
 F=g_ts_\alpha g_s^{-1},\qquad
 \hbar\alpha(c_r-d_r)=0\quad(r\geq2).
\]

Its relative space is then a point and
is otherwise empty.
The proof of
Proposition~\ref{v4-arith:md:prop:arrows}
applies in each component: chain
compatibility makes cohomology preserve
all markings, and the frame action is
free. Duality sends $g$ to
$(g^{-1})^\sharp$ and reverses arrows.
Faces compose the scalars and
degeneracies insert their identities.

These cosets are sets over the fixed
ring $R_2$, not a quotient group or an
exhaustion of arbitrary unmarked native
fields. Recovering a current core and
its constituent moments from intrinsic
field axioms remains a different
recognition problem.

\section{Curvature is a separate comparison}
\label{v4-arith:md:sec:curvature}

Coefficient transposition and a contragredient
module have specified outputs. Neither
identifies two curved field algebras.
That comparison must include changes of
state. A translation class of curvature
can remain invariant under all such changes.

Let $K$ be a Noetherian domain containing
a characteristic-zero field.
A curved vertex superalgebra here has a
degree-one derivation $d$ and a degree-two
state $w$ with
\[
 d^2=w_{(0)},\qquad dw\in TB.
\]
Complete its PBW state space coefficientwise
in an independent degree-zero parameter:
$\widehat B=B_K[[\varepsilon]]$.

Each coefficient is a finite sum of states.
Field identities hold at every finite
$\varepsilon$-precision.
This completion contains $e^{\varepsilon T}a$.
It does not contain an arbitrary derivative
series in the coefficient of
$\varepsilon^0$.

For $\xi\in\widehat B^1$, the changed
differential and curvature are
\begin{equation}\label{v4-arith:md:eq:curved-change}
 d^\xi=d+\xi_{(0)},\qquad
 w^\xi=w+d\xi+\tfrac12\xi_{(0)}\xi.
\end{equation}
A comparison is a continuous unital
graded field isomorphism preserving
$T$ and all modes.
It intertwines the changed differentials
and carries one changed curvature to
the other modulo actual translation.
There is no weight-preservation
requirement.

These changes are well defined at the
chosen completion. Operations are finite
Cauchy sums at each precision, and
\[
 T\widehat B=(TB_K)[[\varepsilon]],\qquad
 \widehat B/T\widehat B
       =(B_K/TB_K)[[\varepsilon]].
\]
Indeed one may choose a translation
primitive separately in each coefficient.
The zero-mode identities give
$(d^\xi)^2=(w^\xi)_{(0)}$.
In the Lie superalgebra
$\widehat B/T\widehat B$, odd Jacobi
then gives $d^\xi w^\xi=0$.
Thus its lift is an actual translation.

\subsection{A detector on a commutative quotient}

Consider
\[
 C_K=\Lambda_K(\theta,T\theta,T^2\theta,\ldots),
 \qquad |\theta|=1,\qquad
 \delta\theta=\theta T\theta.
\]
Its fields are $Y(a,u)b=(e^{uT}a)b$.
All nonnegative modes vanish.
The odd Leibniz rule gives
$\delta^2\theta=0$, so $\delta^2=0$.
Represent a degree-two monomial
$(T^i\theta)(T^j\theta)$, $i<j$, by
the antisymmetric polynomial
$X^iY^j-X^jY^i$.
Translation becomes multiplication by
$X+Y$.

Evaluation at $Y=-X$ has kernel exactly
$(X+Y)$ times the antisymmetric
polynomials. Division by the monic
polynomial $X+Y$ preserves antisymmetry
of the quotient.
Consequently
\begin{equation}\label{v4-arith:md:eq:curvature-detector}
 \ell_3([p])=-\tfrac12[X^3]p(X,-X)
\end{equation}
is well defined on $C_K^2/TC_K^2$.
Its values include
\[
 \ell_3([\theta T^3\theta])=1,\qquad
 \ell_3([T\theta T^2\theta])=-1,\qquad
 \ell_3([\theta T\theta])=0.
\]
The first two are compatible with
$T(\theta T^2\theta)
=T\theta T^2\theta+\theta T^3\theta$.
Extend the detector coefficientwise
to the $\varepsilon$-completion.

\begin{lemma}[Automorphisms of the quotient]
\label{v4-arith:md:lem:quotient-automorphisms}
Every continuous graded
$K[[\varepsilon]]$-linear differential
field automorphism of
$C_K[[\varepsilon]]$ is
$\theta\mapsto e^{sT}\theta$ for
$s\in\varepsilon K[[\varepsilon]]$.
It induces the identity modulo
actual translation.
\end{lemma}

\begin{proof}
The image of $\theta$ is
$P(T)\theta$, where
$P\in K[T][[\varepsilon]]$ is a unit.
Commutation with $\delta$ gives
\[
 P(X)P(Y)(Y-X)=P(X+Y)(Y-X).
\]
Cancel $Y-X$ in the domain.
Invertibility gives $P(0)=1$.
Differentiation at $Y=0$ yields
$P'(X)=sP(X)$, so $P(X)=e^{sX}$.

Its coefficient at $\varepsilon^0$
must be polynomial in $X$.
Characteristic zero and the domain
condition force $s$ to have zero
constant coefficient.
Conversely such an exponential and
its inverse are defined in the
specified completion.
For any completed state $a$,
\[
 (e^{sT}-1)a
   =T\left(\sum_{n\geq1}
           \frac{s^n}{n!}T^{n-1}a\right).
\]
The displayed primitive converges
$\varepsilon$-adically.
\end{proof}

\begin{theorem}[An invariant of changed curvature]
\label{v4-arith:md:thm:curvature}
Suppose a family of curved field
algebras has closed graded field ideals
$I$ preserved by every graded field
isomorphism between its members.
Assume $dI\subset I$ and identify
their completed quotients with
$(C_K[[\varepsilon]],\delta)$.
If
\[
 \pi(w)=\frac c{24}\theta T^3\theta,
\]
then $24\ell_3([\pi(w^\xi)])=c$
for every change of state.
Every comparison above preserves $c$.
\end{theorem}

\begin{proof}
Write $\pi(\xi)=a(T,\varepsilon)\theta$.
Zero modes vanish in the quotient,
so the changed differential remains
$\delta$ there.
The curvature class is
\[
 [\pi(w^\xi)]
  =\frac c{24}[\theta T^3\theta]
      +a(0,\varepsilon)[\theta T\theta].
\]

Every positive power of $T$ in
$\delta\pi(\xi)$ is an actual
translation. The detector kills
the second summand.
A comparison descends to a
differential automorphism of the
quotient. The lemma makes its
action on translation classes the
identity, proving invariance.
\end{proof}

This obstruction is conditional on
the intrinsic quotient and the
specified completion.
A reflection $c\mapsto c'$ can be
implemented by such a comparison
only when $c=c'$.
Enlarging the algebra or changing
the recipient may change the
problem. An uncurving of a larger
algebra therefore requires an
additional comparison with the
original one.

\section{The arithmetic evaluation}
\label{v4-arith:md:sec:mellin}

An arithmetic application needs
a map from its categorical pairing
to an analytic pairing.
We first specify the analytic
reflection and its normalization.
The remaining question can then
be stated as a comparison of
actual maps.

\subsection{Half densities and Mellin reflection}

On $\mathbb R_{>0}$ use Haar measure
$dt/t$. Define
\[
 \mathcal H=L^2(\mathbb R_{>0},
                         t^{1/2}\,dt/t),\qquad
 (\mathcal Jf)(t)=t^{-1/2}f(t^{-1}).
\]
The map
\[
 W:\mathcal H\longrightarrow L^2(\mathbb R,dx),
 \qquad (Wf)(x)=e^{x/4}f(e^x)
\]
is unitary by the substitution $t=e^x$.

The factor $e^{x/4}$ is the square
root of the density relating the
two measures.
Direct calculation gives
$W\mathcal JW^{-1}h(x)=h(-x)$.
Thus
\begin{equation}\label{v4-arith:md:eq:mellin-adjoint}
 \mathcal J^2=1,\qquad
 \mathcal J^\dagger=\mathcal J.
\end{equation}
The adjoint here belongs to the
displayed Hilbert pairing.

For $f\in C_c^\infty(\mathbb R_{>0})$
define
\[
 (\mathcal Mf)(s)=
       \int_0^\infty f(t)t^{s/2}\frac{dt}{t}.
\]
This is entire. Substitution
$t=u^{-1}$ proves
\begin{equation}\label{v4-arith:md:eq:mellin-square}
 \mathcal M(\mathcal Jf)(s)
       =(\mathcal Mf)(1-s).
\end{equation}

Complex conjugation satisfies
$\mathcal M(\overline f)(s)
=\overline{\mathcal Mf(\overline s)}$.
More explicitly,
\[
 \mathcal Mf(1/2+i\tau)
     =\int_{\mathbb R}(Wf)(x)e^{i\tau x/2}\,dx.
\]
These formulas fix both the
reflection and its adjoint.
Using unweighted inversion on the
same Mellin convention would give
$s\mapsto-s$ instead.

\subsection{The Gaussian and both boundary poles}

Fix the Fourier convention
$\widehat f(y)=\int_{\mathbb R}
f(x)e^{-2\pi ixy}\,dx$.
For $f_t(x)=e^{-\pi tx^2}$,
\[
 \widehat f_t(y)=t^{-1/2}e^{-\pi y^2/t}.
\]
For completeness, integration by
parts at $t=1$ gives
$\widehat f_1'(y)=-2\pi y\widehat f_1(y)$.
The Gaussian integral is one,
as follows by squaring it and
using polar coordinates.
The differential equation gives
the formula at $t=1$, and scaling
proves it for every $t>0$.

Periodize $f_t$ over $\mathbb Z$.
The series and all derivatives
converge uniformly on compact
sets. Its Fourier coefficient at
$m$ is $\widehat f_t(m)$.
Both Fourier series converge
absolutely, so evaluation at zero
gives the theta identity
\[
 \theta(t)=\sum_{n\in\mathbb Z}e^{-\pi n^2t}
       =t^{-1/2}\theta(t^{-1}).
\]
This is a lattice theta series.
The oscillator character
$\prod_{n\geq1}(1-q^n)^{-1}$
is a different series: its
coefficient at $q$ is one,
whereas that of
$\sum_{n\in\mathbb Z}q^{n^2}$
is two.

\begin{proposition}[Completed Mellin evaluation]
\label{v4-arith:md:prop:theta-mellin}
For $\operatorname{Re}s>1$,
\[
 \Lambda_\zeta(s)
 :=\pi^{-s/2}\Gamma(s/2)\zeta(s)
 =\frac12\int_0^\infty
       (\theta(t)-1)t^{s/2}\frac{dt}{t}.
\]
Its meromorphic continuation is
\begin{equation}\label{v4-arith:md:eq:polar-mellin}
 \begin{split}
 \Lambda_\zeta(s)
  ={}&-\frac1s-\frac1{1-s}\\
   &+\frac12\int_1^\infty(\theta(t)-1)
       \bigl(t^{s/2}+t^{(1-s)/2}\bigr)\frac{dt}{t}.
 \end{split}
\end{equation}
It has residues $-1$ and $1$ at
zero and one, respectively, and
$\Lambda_\zeta(s)=\Lambda_\zeta(1-s)$.
\end{proposition}

\begin{proof}
Termwise integration is absolutely
convergent for $\operatorname{Re}s>1$.
Substitution $u=\pi n^2t$ in each
positive-integer term gives
$\pi^{-s/2}\Gamma(s/2)n^{-s}$.
The two signs of $n$ cancel the
factor $1/2$.

Split the integral at one.
On $(0,1)$ use
\[
 \theta(t)-1
  =t^{-1/2}(\theta(t^{-1})-1)
       +t^{-1/2}-1.
\]
Inverting $t$ in the first summand
gives the integral with exponent
$(1-s)/2$ over $(1,\infty)$.
The last two terms integrate to
$1/(s-1)-1/s$.
The remaining integral is entire:
the theta tail decays exponentially,
uniformly against every power
and every power of $\log t$
on compact $s$-sets.
The formula proves the residues
and reflection.
\end{proof}

The theta tail is outside
$\mathcal H$: near zero it has
leading term $t^{-1/2}$.
Its Mellin evaluation therefore
uses the displayed subtraction
and continuation.
It cannot be obtained by applying
an everywhere-defined Hilbert
pairing to that state.

On boundary monomials,
$\mathcal J1=t^{-1/2}$ and
$\mathcal J(t^{-1/2})=1$.
Their integrals over $(0,1)$
produce the two polar terms.
An arithmetic comparison must
retain these boundary contributions
as well as its native vacuum.
No identification of that vacuum
with either boundary monomial
has been assumed.

The entire normalization
$\xi(s)=s(s-1)\Lambda_\zeta(s)/2$
removes the poles and is still
reflection invariant.
Multiplication by this polynomial
is not a construction of the
missing arithmetic evaluation.
Likewise Euler inversion to
$1/\zeta(s)$ changes the scalar
function without constructing
the involution in
\eqref{v4-arith:md:eq:mellin-square}.

\subsection{Detecting a central defect}

Let $V$ be a complex vector space,
$J:V\to V$ an involution, and
$E:V\to\mathscr M(\mathbb C)$ a
linear map to meromorphic functions.
Fix $v,w\in V$ and scalars
$\kappa,\kappa^!$.
The following criterion keeps the
categorical construction separate
from its consequence.

\begin{proposition}[Detection of a specified central defect]
\label{v4-arith:md:prop:central-defect}
Suppose
\[
 v-Jv=(\kappa+\kappa^!)w,\qquad
 E(v)(s)=\Lambda_\zeta(s),\qquad
 E(Jv)(s)=\Lambda_\zeta(1-s),
\]
and $E(w)$ is not identically zero.
Then $\kappa+\kappa^!=0$ if and
only if
$\Lambda_\zeta(s)=\Lambda_\zeta(1-s)$.
The same equivalence holds for
the entire normalization $\xi$.
\end{proposition}

\begin{proof}
Apply $E$ to the defect equation:
\[
 \Lambda_\zeta(s)-\Lambda_\zeta(1-s)
       =(\kappa+\kappa^!)E(w)(s).
\]
A scalar multiple of a nonzero
meromorphic function vanishes
exactly when the scalar does.
The polynomial $s(s-1)/2$ is
nonzero and reflection invariant,
which proves the assertion for
$\xi$.
\end{proof}

For the forward implication,
equivariance
$E(Jx)(s)=E(x)(1-s)$ and the
defect equation suffice.
For the converse, detection of
the stated defect is necessary.
For example, on $\mathbb C^2$
let $J(x,y)=(y,x)$ and
$E(x,y)=x+y$.
The evaluations of $v=(1,0)$
and $Jv$ agree, but
$v-Jv=(1,-1)\ne0$ lies in
$\ker E$.
Full injectivity of $E$ is
unnecessary. Nonvanishing on
the specified defect line
is the relevant hypothesis.

The scalar reflection has been
proved independently above.
Consequently this proposition
tests a proposed arithmetic
construction: if its defect and
evaluation identities hold,
its level sum must vanish.
An equivalence of categories
without those identities gives
no value for that sum.

\subsection{Local assembly and the remaining comparison}

Suppose presentable stable
$R$-linear categories
$\mathcal C_v,\mathcal D_v$
have chosen objects $c_v,d_v$
at each place.
For a finite set $F$ take
their tensor products over $R$.
Use insertion of the chosen
objects as transition functors.
Compatible local equivalences
$\Psi_v:\mathcal C_v\simeq\mathcal D_v$
with coherent identifications
$\Psi_v(c_v)\simeq d_v$ give
an equivalence
\[
 \varinjlim_F\bigotimes_{v\in F,R}\mathcal C_v
 \simeq
 \varinjlim_F\bigotimes_{v\in F,R}\mathcal D_v.
\]
Indeed they give an equivalence
at every finite tensor stage
and commute with its transition
maps. Colimits preserve this
equivalence of diagrams.

For arithmetic categories one
must construct the common
coefficient action, local
functors and compatible
pointings. A unital map
$\mathbb F_p\to\mathbb C$
cannot supply a coefficient
comparison, since it would
send $p\cdot1=0$ to
$p\cdot1\ne0$.
A common integral base with
two derived fibers is a
different, possible construction.
Identifying the categorical
colimit with an analytic
restricted tensor product
also requires its own map.

\begin{question}[Completed arithmetic duality and evaluation]
\label{v4-arith:md:qn:arithmetic}
Construct the archimedean
coefficient complex and the
pairing of a functorial completed
arithmetic Koszul duality.
Construct an operation-preserving
comparison to theta and Mellin
reflection with its specified
half-density normalization and
adjoint.
The comparison must transport
both polar terms and the native
vacuum, as well as the products,
collision operations and
completions on its declared
domain.
\end{question}

The first missing map is from
this arithmetic coefficient
complex and its duality pairing
to the analytic recipient with
the evaluation above.
To obtain a level criterion,
it must also supply the defect
identity and detect its defect
line.
The coefficient anti-equivalence,
the marked native examples and
the scalar theta theorem each
provide positive structures.
Their comparison remains the
mathematical construction asked
for in the question.

\section*{References for the duality theorems}

Pontryagin duality is treated by David W.~Roeder,
\emph{Category theory applied to Pontryagin duality},
Pacific Journal of Mathematics \textbf{52} (1974), 519--527.
Matlis duality and derived completion use the Stacks Project
tags specified in their proofs.

The field-coefficient correspondence and its naturality use
Leonid Positselski,
\emph{Two kinds of derived categories, Koszul duality,
and comodule--contramodule correspondence}, arXiv:0905.2621.
Sections 4.4, 5.2 and 5.4 give the resolution models,
the correspondence and its derived-adjoint comparison.
Sections 6.3, 6.4 and 6.9 give the twisting-cochain comparisons.

The contragredient formula is Yi-Zhi Huang,
\emph{Introduction to vertex operator algebras, III},
arXiv:q-alg/9504019, equation (1.26).
The formal invariance theorem is Ezra Getzler,
\emph{Lie theory for nilpotent \(L_\infty\)-algebras},
Annals of Mathematics \textbf{170} (2009), 271--301, Theorem 4.8.

Bruce Ferrero and Lawrence C.~Washington,
\emph{The Iwasawa invariant \(\mu_p\) vanishes for abelian number fields},
Annals of Mathematics \textbf{109} (1979), 377--395.
